% =====================================================================
% main.tex -- HSC Physics ESSENTIALS, Module 6: Electromagnetism.
% The condensed, knowledge-only companion to the full revision guide
% (src/mod6/main.tex). Same document class and design system.
%
% Build with:
%     ./build.sh mod6_essentials
%     (final copy: build_book.sh mod6_essentials Module6_Electromagnetism_Essentials.pdf)
%
% Book order (Essentials editions):
%   cover - How to use these Essentials (below) - Contents
%   chapters ch1.tex ... ch5.tex
%   All formulas at a glance                                      (formulas.tex)
%   Key experiments at a glance                                   (experiments.tex)
%   Glossary                                                      (glossary.tex)
% No questions, answers, practice set or "Answering HSC questions" part.
% =====================================================================
\documentclass[11pt,a4paper,oneside,openany]{book}
\makeatletter\def\input@path{{./}{../common/}}\makeatother
\usepackage{hscphys}
\IfFileExists{./module_info.tex}{% module_info.tex -- Module 6 ESSENTIALS edition (condensed, knowledge only, no questions).
% Final PDF: build_book.sh mod6_essentials Module6_Electromagnetism_Essentials.pdf  ->  output/Module6_Electromagnetism_Essentials.pdf
\def\modulenumber{6}
\def\moduletitle{Electromagnetism}
\makeatletter
\renewcommand\hsccoversubtitle{HSC Physics Essentials — key ideas and formulas}
\renewcommand\hsccoverchips{%
  \hsc@coverchip{hscbigidea}{\faLightbulb}{Key ideas}\hspace{2mm}%
  \hsc@coverchip{hscformula}{\faIcon{square-root-alt}}{Formulas}\hspace{2mm}%
  \hsc@coverchip{hscexp}{\faFlask}{Key experiments}\hspace{2mm}%
  \hsc@coverchip{hscbigidea}{\faIcon{book}}{Glossary}}
\renewcommand\hsccoverfooter{Condensed edition \textperiodcentered\ knowledge points only \textperiodcentered\ no questions}
\makeatother
}{}
% PDF title of this edition (hscphys.sty sets the full guide's title; this runs after it)
\AtBeginDocument{\hypersetup{
  pdftitle={Module \modulenumber: \moduletitle\space - HSC Physics Essentials}}}

\begin{document}

%% ---------------------------------------------------------------- front
\makecover

\unnumberedchapter[\faInfoCircle]{How to use these Essentials}

This is the short, knowledge-only companion to the full \emph{HSC Physics Revision
Guide} for Module~\modulenumber\ (\moduletitle). It covers every syllabus dot point (NSW
Physics Stage~6 Syllabus, 2017). It has no questions: the worked examples, Quick checks
and practice set with answers are in the full guide.

\subsection*{What the coloured boxes mean}

\begin{sidebyside}[0.5]
  \begin{bigidea}[before skip=0pt, after skip=0pt, equal height group=esslegA, fontupper=\normalsize]
    The main point of a section in one to three sentences. Read it first.
  \end{bigidea}
\nextside
  \begin{keywords}[before skip=0pt, after skip=0pt, equal height group=esslegA]
    \kw{Term}{the meaning in plain words. Other new words are in \term{bold}.}
  \end{keywords}
\end{sidebyside}

\begin{formulacard}[unbreakable, before skip=6pt, after skip=6pt]{Name of the formula}
  \cardformula{V = IR}
  The symbols with their SI units, when to use it and a common trap. The badge is
  about the official HSC formula sheet:
  \begin{itemize}[itemsep=0pt, topsep=1pt]
    \item \hscbadgeonsheet\ \ printed on the sheet;
    \item \hscbadgenotonsheet\ \ learn it, or learn how to get it;
    % the combine badge is not an on-line box: lower it so its text sits on the baseline
    \item \raisebox{-2.3pt}{\hscbadgecombine{two formulas}}\ \ join two formulas.
  \end{itemize}
\end{formulacard}

\begin{sidebyside}[0.5]
  \begin{experiment}[unbreakable, before skip=0pt, after skip=0pt, equal height group=esslegB]{}
    A practical or historical experiment: aim, method, results, conclusion.
  \end{experiment}
\nextside
  \begin{summarybox}[unbreakable, before skip=0pt, after skip=0pt, equal height group=esslegB]{Remember}
    \emph{In this chapter} (start) and \emph{Remember} (end): the must-know points.
  \end{summarybox}
\end{sidebyside}

\subsection*{How to use this book}
\begin{itemize}[itemsep=1pt]
  \item \textbf{Read one chapter a day.} Start each section with its key idea. Learn the
        formulas with the red badge.
  \item If a part is not clear, read the same section in the full guide.
  \item \textbf{The night before the exam} (Wednesday 4~November 2026), re-read the
        \emph{Remember} boxes and \emph{All formulas at a glance}.
\end{itemize}

% Local layout fix (common/ is shared, so it is made here): the contents of this short
% book fit on one page with a little less space between the entries. The setting is
% written into the .toc file, so it changes only the entries (the contents file is read
% inside a group), not the page heading.
\addtocontents{toc}{\protect\setlength{\protect\parskip}{0pt plus 1pt}\protect\setstretch{1.0}}
\tableofcontents

%% ---------------------------------------------------------------- chapters
\IfFileExists{./ch1.tex}{\setcounter{chapter}{0}% ch1.tex -- Module 6 ESSENTIALS, Chapter 1: Charged particles in electric and magnetic fields
% Condensed from the final full book: src/mod6/ch1.tex (same notation, same facts).
% Knowledge points only: no worked examples, no questions, no answers.
% The few numbers kept (2e12, 1840, 2e/m_e) are checked by
\chapter{Charged Particles in Electric and Magnetic Fields}\label{ch1:chap}
\covers{M6.IQ1.01, M6.IQ1.02, M6.IQ1.03, M6.IQ1.04}

\begin{summarybox}{In this chapter}
\begin{itemize}
  \item Charges between parallel plates: the field, the force, the work and the speed.
  \item Curved (parabolic) paths between plates, compared with projectiles.
  \item The magnetic force on a moving charge, and circular paths.
  \item Electric, magnetic and gravitational fields compared.
\end{itemize}
\end{summarybox}

%% ==================================================================
\section{Charges in a uniform electric field}\label{ch1:sec:efield}
\covers{M6.IQ1.01}

\begin{bigidea}
Two charged parallel plates make a uniform electric field. A charge between
them feels a constant force, so it has a constant acceleration. The field does
work $W = qV$ on the charge, and this becomes kinetic energy.
\end{bigidea}

\begin{keywords}
  \kw{Electric field strength ($E$)}{the force per coulomb. Unit: \si{V.m^{-1}} $=$ \si{N.C^{-1}}.}
  \kw{Uniform field}{a field with the same strength and direction everywhere.}
  \kw{Electron volt (eV)}{the energy one electron (or proton) gains through
    \qty{1}{V}. $\qty{1}{eV} = \qty{1.602e-19}{J}$. Energy in eV $=$ (charge in
    units of $e$) $\times$ (voltage in V).}
\end{keywords}

\begin{sidebyside}[0.5]
  \centering
  \begin{tikzpicture}[x=0.74cm, y=0.8cm]
    \draw[field] (0.18,2.0) .. controls (-0.15,1.45) and (-0.15,0.55) .. (0.18,0);
    \draw[field] (5.82,2.0) .. controls (6.15,1.45) and (6.15,0.55) .. (5.82,0);
    \foreach \x in {0.75,1.5,...,5.25} \draw[field] (\x,2.0) -- (\x,0);
    \fill[hscink!75] (0,2.0) rectangle (6,2.17);
    \fill[hscink!75] (0,-0.17) rectangle (6,0);
    \foreach \x in {0.375,1.125,...,5.7} {
      \node[font=\small\bfseries] at (\x,2.55) {$+$};
      \node[font=\small\bfseries] at (\x,-0.55) {$-$};
    }
    \draw[line width=0.8pt] (6,2.08) -- (7.2,2.08) -- (7.2,1.15);
    \draw[line width=0.8pt] (6.9,1.15) -- (7.5,1.15);
    \draw[line width=2.2pt] (7.05,0.9) -- (7.35,0.9);
    \draw[line width=0.8pt] (7.2,0.9) -- (7.2,-0.08) -- (6,-0.08);
    \node[diagramlabel, right] at (7.5,1.03) {$V$};
    \draw[<->, line width=0.6pt, hscmuted] (-0.8,0) -- (-0.8,2.0)
      node[midway, left, text=hscink, font=\small] {$d$};
    \node[font=\small, text=hscfield, fill=white, inner sep=1.5pt] at (3.375,1.0) {$E$};
    \draw[force] (1.875,1.6) -- (1.875,0.5);
    \fill[object] (1.875,1.6) circle (4.5pt);
    \node[font=\scriptsize\bfseries] at (1.875,1.6) {$+$};
    \node[diagramlabel, font=\footnotesize, text=hscforce, inner sep=0.5pt] at (1.875,0.25) {$F$};
    \draw[force] (4.875,0.4) -- (4.875,1.5);
    \fill[object] (4.875,0.4) circle (4.5pt);
    \node[font=\scriptsize\bfseries] at (4.875,0.4) {$-$};
    \node[diagramlabel, font=\footnotesize, text=hscforce, inner sep=0.5pt] at (4.875,1.75) {$F$};
  \end{tikzpicture}
  \captionof{figure}{Plates joined to a battery (voltage $V$). A positive
    charge is pushed along the field; an electron is pushed the other
    way.}\label{ch1:fig:plates}
\nextside
  \begin{itemize}
    \item An \term{electric field} is a region where a charge feels a force. Its
          direction is the force on a \emph{positive} charge.
    \item \term{Field lines} go from $+$ to $-$. Between the plates
          (\cref{ch1:fig:plates}) they are straight, parallel and evenly spaced.
    \item Near the ends they bulge out (the \term{edge effect}). Calculations
          ignore it.
  \end{itemize}
\end{sidebyside}

\begin{itemize}
  \item The \term{elementary charge} is $e = \qty{1.602e-19}{C}$ (the sheet prints
        $q_\mathrm{e}$ as negative). Use the size; find the direction separately.
  \item A positive charge is pushed along the field; an electron is pushed against
        it. Released from rest, a charge moves in a straight line, speeding up
        steadily.
  \item Constant force, so constant \term{acceleration}:
        $a = \frac{qE}{m} = \frac{qV}{md}$ (combine $\vec{F} = q\vec{E}$,
        $\vec{F}_\mathrm{net} = m\vec{a}$, $E = \frac{V}{d}$). A bigger
        \term{charge-to-mass ratio} $\frac{q}{m}$ gives a bigger $a$.
  \item Ignore gravity for electrons, protons and \term{ions} (charged atoms): in
        only \qty{100}{V.m^{-1}}, the electric force on an electron is about
        \num{2e12} times its weight.
  \item An oil drop held still: electric force up = weight down, $qE = mg$ (not on
        the sheet). Its charge is $q = ne$, $n$ a whole number (Millikan, Module 8).
  \item The field does \term{work} (transfers energy): the charge loses
        \term{electric potential energy} and gains kinetic energy, like a ball
        rolling downhill. $V = \frac{\Delta U}{q}$ (on the sheet), so $W = qV$.
        The gap does not matter: closer plates push harder, but for a shorter
        distance.
\end{itemize}

\begin{formulacard}{Uniform electric field: field, force and work}
  \cardformula{E = \frac{V}{d} \qquad \vec{F} = q\vec{E} \qquad W = qV \qquad W = qEd}
  \begin{symbols}
    \sym{$E$, $V$}{electric field strength; potential difference (voltage)}{\si{V.m^{-1}}, \si{V}}
    \sym{$d$}{plate gap ($E = \frac{V}{d}$); distance moved along $E$ ($W = qEd$)}{\si{m}}
    \sym{$\vec{F}$, $q$}{electric force on the charge; the charge (with its sign in
      $\vec{F} = q\vec{E}$, its size in $W = qV$)}{\si{N}, \si{C}}
    \sym{$W$}{work done on the charge by the field}{\si{J}}
  \end{symbols}
  \whentouse{$E = \frac{V}{d}$, $W = qEd$: uniform field only.
    $\vec{F} = q\vec{E}$, $W = qV$: any field.}
  \derive{across the whole gap, $W = Fd = qEd = q\left(\dfrac{V}{d}\right)d = qV$.}
  \trap{change mm, cm and kV to m and V first. For a negative charge $\vec{F}$ points
    against $\vec{E}$.}
\end{formulacard}

\begin{formulacard}[combine={$W = qV$ and $K = \frac{1}{2}mv^2$}]{Speed gained from rest}
  \cardformula{\frac{1}{2}mv^2 = qV \quad\text{so}\quad v = \sqrt{\frac{2qV}{m}}}
  \begin{symbols}
    \sym{$v$}{final speed (starting from rest)}{\si{m.s^{-1}}}
    \sym{$m$, $q$}{mass; size of the charge}{\si{kg}, \si{C}}
    \sym{$V$}{accelerating voltage}{\si{V}}
  \end{symbols}
  \whentouse{a charge sped up from rest through $V$ ($v$ much less than $c$).}
  \derive{work done = gain in kinetic energy: $qV = \frac{1}{2}mv^2 - 0$.}
  \trap{do not forget the 2 and the square root. A charge already moving at $u$:
    $qV = \frac{1}{2}mv^2 - \frac{1}{2}mu^2$.}
\end{formulacard}

An \term{electron gun} uses this. A hot \term{cathode} (negative plate) gives off
electrons. The field speeds them up, and they leave through a hole in the
positive plate as a \term{beam}.

%% ==================================================================
\section{Curved paths in an electric field}\label{ch1:sec:paths}
\covers{M6.IQ1.02}

\begin{bigidea}
A charge that enters moving parallel to the plates follows a parabola, like a
ball thrown horizontally. Its velocity along the plates stays the same; across
the plates it has a constant acceleration.
\end{bigidea}

\begin{sidebyside}[0.56]
  \centering
  \begin{tikzpicture}[x=0.7cm, y=0.7cm]
    \foreach \x in {0.6,1.4,...,5.4} \draw[line width=0.8pt, hscfield!45,
      postaction={decorate}, decoration={markings, mark=at position 0.78 with
      {\arrow{Stealth[length=2.2mm, width=1.7mm]}}}] (\x,1.5) -- (\x,-1.5);
    \draw[line width=0.4pt, hscmuted] (0,1.5) -- (-2.15,1.5);
    \draw[line width=0.4pt, hscmuted] (0,-1.5) -- (-2.15,-1.5);
    \draw[<->, line width=0.6pt, hscmuted] (-2.0,-1.5) -- (-2.0,1.5)
      node[pos=0.25, left, text=hscink, font=\small] {$d$};
    \fill[hscink!75] (0,1.5) rectangle (6,1.66);
    \fill[hscink!75] (0,-1.66) rectangle (6,-1.5);
    \node[diagramlabel, above] at (3,1.66) {positive plate};
    \node[diagramlabel, below] at (1.6,-1.66) {negative plate};
    \draw[guide] (0,0) -- (8.3,0);
    \draw[hscink, line width=1pt, domain=0:6, samples=40] plot (\x, {1.1*\x*\x/36});
    \draw[hscink, line width=1pt, dashed] (7.8,1.76) -- (8.6,2.053);
    \fill[object] (-0.15,0) circle (4.5pt);
    \node[font=\scriptsize\bfseries] at (-0.15,0) {$-$};
    \draw[velocity] (-1.6,0) -- (-0.45,0) node[pos=0.4, above, text=hscvel] {$u$};
    \draw[force] (3,0.275) -- (3,1.15) node[left, text=hscforce] {$F$};
    \fill[hscink] (3,0.275) circle (1.8pt);
    \draw[velocity, hscvel!60] (6,1.1) -- (7.8,1.1) node[below, text=hscvel] {$v_x = u$};
    \draw[velocity, hscvel!60] (7.8,1.1) -- (7.8,1.76) node[pos=0.5, right, text=hscvel] {$v_y$};
    \draw[velocity] (6,1.1) -- (7.8,1.76) node[pos=0.55, above, text=hscvel] {$v$};
    \draw[<->, line width=0.6pt, hscmuted] (5.8,0) -- (5.8,1.06) node[midway, right, text=hscink, font=\small] {$y$};
    \draw[<->, line width=0.6pt, hscmuted] (0,-2.9) -- (6,-2.9) node[pos=0.75, fill=white, text=hscink, font=\small] {$L$};
    \node[font=\small, text=hscfield, fill=white, inner sep=1pt] at (1.0,-0.8) {$E$};
  \end{tikzpicture}
  \captionof{figure}{An electron curves towards the positive plate in a
    parabola, then leaves in a straight line.}\label{ch1:fig:parabola}
\nextside
  \begin{itemize}
    \item The \term{trajectory} (path) in \cref{ch1:fig:parabola}: along the
          plates there is no force, so $v_x = u$ is constant and $x = ut$.
    \item Across the plates $a = \dfrac{qE}{m}$ is constant. With $u_y = 0$,
          $y = \frac{1}{2}at^2$.
    \item The path bends towards the plate that attracts the charge.
    \item After the plates there is no force: a straight line.
  \end{itemize}
\end{sidebyside}

\begin{itemize}
  \item \textbf{Why a \term{parabola}:} put $t = \frac{x}{u}$ into
        $y = \frac{1}{2}at^2$: $y = \left(\dfrac{qE}{2mu^2}\right)x^2$, so
        $y \propto x^2$. A thrown ball has the same shape, with $g$ for $\frac{qE}{m}$.
  \item The \term{deflection} (sideways push) at the end of plates of length $L$:
        $t = \frac{L}{u}$, $y = \dfrac{qEL^2}{2mu^2}$ (combine
        $s = ut + \frac{1}{2}at^2$ and $a = \frac{qE}{m}$). $L$ is the plate
        length, not the gap. A faster particle is deflected less.
  \item Exit velocity: $v_y = at$, $v = \sqrt{v_x^2 + v_y^2}$,
        $\theta = \tan^{-1}\!\left(\frac{v_y}{v_x}\right)$. Entering at an angle:
        split $u$ into $u\cos\theta$ and $u\sin\theta$ first.
  \item \textbf{Like a projectile:} a parabola (constant force in one direction,
        none at right angles), and the same equations in each direction.
  \item \textbf{Unlike a projectile:} $a = \frac{qE}{m}$ depends on charge and mass,
        whereas $g$ is the same for every mass. The force can point either way
        (sign of $q$), and $a$ is huge for electrons.
\end{itemize}

%% ==================================================================
\section{Charges in a uniform magnetic field}\label{ch1:sec:bfield}
\covers{M6.IQ1.03, M6.IQ1.04}

\begin{bigidea}
A magnetic field pushes only on a \emph{moving} charge, at right angles to both
the velocity and the field. The push changes the direction of motion but never
the speed. So a charge moving at \ang{90} to a uniform field moves in a circle.
\end{bigidea}

\subsection{The magnetic force}
\begin{itemize}
  \item A \term{magnetic field} ($B$) is a region where a magnet or a moving
        charge feels a magnetic force. Unit: \term{tesla} (T).
  \item A charge at rest feels no force. Only $v_\perp = v\sin\theta$ counts: the
        part of $v$ \term{perpendicular} (at \ang{90}) to $B$. Along $B$: no force.
  \item The force is perpendicular to \emph{both} $v$ and $B$.
        $\bullet\,\bullet\,\bullet$ = field out of the page;
        $\times\,\times\,\times$ = field into the page.
\end{itemize}

\begin{formulacard}{Magnetic force on a moving charge}
  \cardformula{F = qv_\perp B = qvB\sin\theta}
  \begin{symbols}
    \sym{$F$}{magnetic force on the charge}{\si{N}}
    \sym{$q$, $v$, $v_\perp$}{size of the charge; speed; part of $v$ perpendicular to $B$}{\si{C}, \si{m.s^{-1}}}
    \sym{$B$, $\theta$}{field strength; angle between $v$ and $B$}{\si{T}, \si{\degree}}
  \end{symbols}
  \whentouse{a charge moving through a magnetic field.}
  \trap{$\theta$ is measured from the \emph{field} (use $\sin\theta$). $v$ in the page
    and $B$ into or out of the page: $\theta = \ang{90}$.}
\end{formulacard}

\begin{sidebyside}[0.5]
  \centering
  \begin{tikzpicture}[x=0.85cm, y=0.85cm]
    \foreach \x/\h in {0.05/0.95, 0.47/1.22, 0.89/1.38, 1.31/1.2}
      \draw[line width=0.8pt, fill=hscbrand!7, rounded corners=5pt] (\x,1.35) rectangle ++(0.4,\h);
    \draw[line width=0.8pt, fill=hscbrand!7, rounded corners=5pt] (1.4,0.38) rectangle (2.75,0.82);
    \draw[line width=0.8pt, fill=hscbrand!7, rounded corners=7pt] (0,0) rectangle (1.75,1.7);
    \draw[line width=0.8pt] (0.25,0) -- (0.25,-0.4);
    \draw[line width=0.8pt] (1.5,0) -- (1.5,-0.4);
    \draw[->, line width=1.3pt, hscfield] (-0.5,0.6) -- (-0.5,2.7) node[above, font=\small] {$B$};
    \node[diagramlabel, left, align=right] at (-0.65,1.7) {fingers:\\ field};
    \draw[velocity] (2.9,0.6) -- (3.9,0.6) node[right, text=hscvel] {$v$};
    \node[diagramlabel, align=left, anchor=north west] at (2.85,0.35) {thumb: velocity\\ of a $+$ charge};
    \draw[hscforce, line width=1.1pt] (0.87,0.8) circle (0.2);
    \fill[hscforce] (0.87,0.8) circle (0.06);
    \node[diagramlabel, text=hscforce, align=center, below] at (0.87,-0.45) {palm faces you:\\ $F$ out of the page};
  \end{tikzpicture}
  \captionof{figure}{The right-hand palm rule.}\label{ch1:fig:hand}
\nextside
  \textbf{Direction:} the \term{right-hand palm rule} (\cref{ch1:fig:hand}).
  Flat \emph{right} hand: fingers along $B$, thumb along $v$ of a
  \emph{positive} charge. The palm faces the force. For a negative charge,
  reverse it. So with $B$ into the page and $v$ to the right, the force is up
  the page on a $+$ charge, and down on an electron.
\end{sidebyside}

\subsection{Circular paths}
The acceleration $a = \frac{F}{m}$ is along $F$: perpendicular to the field.
The path is a circle (\cref{ch1:fig:circles}) because:
\begin{enumerate}
  \item $F$ is always at \ang{90} to $v$, so $W = Fs\cos\ang{90} = 0$: no work is done.
  \item No work, so the speed is constant and $F = qvB$ keeps the same size.
  \item A constant-size force always at \ang{90} to $v$ is a \term{centripetal
        force} (net force towards the centre): \term{uniform circular motion}
        (moving in a circle at a constant speed).
\end{enumerate}

\begin{figure}
  \centering
  \begin{tikzpicture}[x=0.74cm, y=0.74cm]
    % (a) positive charge, (b) electron; B into the page
    \foreach \cx/\sgn/\chg/\lab in {0/1/+/{(a) $+$ charge:\\ anticlockwise},
                               4.2/-1/-/{(b) electron:\\ clockwise}} {
      \begin{scope}[xshift=\cx*0.74 cm]
        \foreach \x/\y in {-0.5/0.5, -0.5/-0.5,
            1.6/1.6, -1.6/1.6, 1.75/-0.95, -1.75/-0.95}
          \node[text=hscfield!40, font=\small] at (\x,\y) {$\times$};
        \draw[hscink!55, line width=0.8pt] (0,0) circle (1.2);
        \fill[hscink] (0,0) circle (1.3pt);
        \foreach \a in {0,90,180,270} {
          \coordinate (P) at (\a:1.2);
          \draw[velocity] (P) -- ++({\a+\sgn*90}:0.75);
          \draw[force] (P) -- ($(P)!0.5!(0,0)$);
          \fill[object] (P) circle (4.2pt);
          \node[font=\scriptsize\bfseries] at (P) {$\chg$};
        }
        \node[diagramlabel, text=hscvel] at (1.55,{\sgn*0.55}) {$v$};
        \node[diagramlabel, text=hscforce] at (0.82,{-\sgn*0.26}) {$F$};
        \node[diagramlabel, align=center] at (0,-2.1) {\lab};
      \end{scope}
    }
    % (c) half a circle
    \begin{scope}[xshift=9.9*0.74cm, yshift=0.35*0.74cm]
      \foreach \x/\y in {-2.4/0.4, -2.4/1.3, 0.6/0.55, 0.6/1.35,
          -1.5/0.3, -1.0/1.35}
        \node[text=hscfield!40, font=\small] at (\x,\y) {$\times$};
      \draw[line width=0.7pt, hscmuted] (-2.75,0) -- (0.95,0)
        node[right, text=hscmuted, font=\small, align=left] {edge of\\ field};
      \draw[hscink, line width=1pt] (0,-1.0) -- (0,0);
      \draw[hscink, line width=1pt] (0,0) arc (0:180:1.0);
      \draw[hscink, line width=1pt, ->] (-2.0,0) -- (-2.0,-1.0);
      \draw[velocity] (0,-1.0) -- (0,-0.15) node[pos=0.35, right, text=hscvel] {$v$};
      \fill[object] (0,-1.0) circle (4.2pt);
      \node[font=\scriptsize\bfseries] at (0,-1.0) {$+$};
      \coordinate (Q) at ($(-1.0,0)+(45:1.0)$);
      \draw[force] (Q) -- ($(Q)!0.55!(-1.0,0)$);
      \node[diagramlabel, text=hscforce] at ($(Q)+(-0.6,-0.12)$) {$F$};
      \fill[hscink] (-1.0,0) circle (1.3pt);
      \draw[<->, line width=0.6pt, hscmuted] (-2.0,-1.45) -- (0,-1.45)
        node[midway, fill=white, text=hscink, font=\small, inner sep=1pt] {$2r$};
      \node[diagramlabel, align=center] at (-1.0,-2.45) {(c) half a circle,\\ out $2r$ away};
    \end{scope}
    % (d) helix
    \begin{scope}[xshift=12.9*0.74cm]
      \draw[field] (0,-1.15) -- (3.4,-1.15) node[right, text=hscfield] {$B$};
      \draw[hscink, line width=0.9pt, domain=0:1440, samples=200, smooth]
        plot ({0.25 + \x/480 + 0.3*sin(\x)}, {0.8*cos(\x)});
      \node[diagramlabel, align=center] at (1.7,-2.1) {(d) $v$ at an angle\\ to $B$: a helix};
    \end{scope}
  \end{tikzpicture}
  \caption{Paths in a uniform magnetic field ($B$ into the page in (a)--(c)).
    Blue: velocity; red: force, towards the centre.}\label{ch1:fig:circles}
\end{figure}

\begin{formulacard}[notonsheet]{Radius and period in a magnetic field}
  \cardformula{r = \frac{mv}{qB} \qquad\quad T = \frac{2\pi m}{qB}}
  \begin{symbols}
    \sym{$r$, $T$}{radius of the circle; period (time for one turn)}{\si{m}, \si{s}}
    \sym{$m$, $q$}{mass; size of the charge}{\si{kg}, \si{C}}
    \sym{$v$, $B$}{speed (perpendicular to $B$); field strength}{\si{m.s^{-1}}, \si{T}}
  \end{symbols}
  \whentouse{a charge at \ang{90} to a uniform $B$, with no other force.}
  \derive{magnetic force = centripetal force (both on the sheet):
    $qvB = \dfrac{mv^2}{r}$, so $r = \dfrac{mv}{qB}$. Then
    $T = \dfrac{2\pi r}{v} = \dfrac{2\pi m}{qB}$ (from $v = \frac{2\pi r}{T}$).}
  \trap{neither formula is on the sheet. Start from $qvB = \dfrac{mv^2}{r}$.}
\end{formulacard}

\begin{itemize}
  \item Heavier or faster: a wider circle ($r \propto v$). More charge or a
        stronger field: a tighter circle. $T$ does not depend on the speed.
  \item Crossing the edge of a large field at \ang{90}: half a circle, out
        through the same edge $2r$ away. A \term{mass spectrometer} uses this to
        sort ions: with the same charge and speed, heavier ions land further away.
  \item $v$ at another angle to $B$: $v_\perp$ makes a circle, $v_\parallel$
        (no force) carries it along the field. The path is a \term{helix} (a
        spiral, like a stretched spring).
\end{itemize}

\subsection{Comparing electric, magnetic and gravitational fields}
\covers{M6.IQ1.04}

In gravity and in an electric field the force has a fixed direction: a
parabola. In a magnetic field the force turns with the velocity: a circle
(\cref{ch1:tab:EB}).

\begin{table}
  \caption{A charge in a uniform electric field compared with a charge in a
    uniform magnetic field.}\label{ch1:tab:EB}
  \centering\small
  \begin{tabular}{@{}L{34mm}L{54mm}L{52mm}@{}}
    \toprule
    & \colhead{Uniform electric field} & \colhead{Uniform magnetic field}\\
    \midrule
    Charge at rest & feels a force $qE$ & feels no force\\
    Direction of force & along the field ($+$) or against it ($-$) &
      perpendicular to both $v$ and $B$\\
    Moving along the field & speeds up or slows down in a line & no force:
      constant velocity\\
    Entering at \ang{90} & parabola & circle (or an arc)\\
    Acceleration & constant size and direction & constant size; turns (centripetal)\\
    Work and speed & work $qV$ is done; speed changes & no work; speed constant\\
    \bottomrule
  \end{tabular}
\end{table}

\begin{itemize}
  \item \textbf{Like all uniform circular motion} (tension on a ball on a string,
        friction on a car, gravity on a satellite): $F$ is at \ang{90} to $v$, so
        there is no work and the speed is constant.
  \item \textbf{Differences:} $qvB$ exists only because the charge moves. For a
        charge $r \propto v$ and $T$ does not depend on speed (why a cyclotron, a
        machine that speeds up particles in circles, works). A satellite in a bigger orbit is slower and takes longer.
  \item \textbf{Crossed fields} ($E$ and $B$ at right angles): when $qE = qvB$ the
        forces cancel. Only the speed $v = \frac{E}{B}$ (not on the sheet) goes
        straight through: a \term{velocity selector} (Thomson, Module 8).
\end{itemize}

%% ==================================================================
\section{Key experiments}\label{ch1:sec:experiments}
\covers{M6.IQ1.01, M6.IQ1.02}

\begin{experiment}{Simulating charges between parallel plates}
\aim{To test the path of a charge between the plates, and the speed gained from rest.}
\method{A computer simulation. (1)~An electron enters between the plates (\cref{ch1:fig:parabola}).
  Change only $V$ and record $y$; repeat with half the speed, and with a proton;
  record $y$ at several $x$. (2)~An electron gun: record the exit speed $v$ for
  several $V$, then double the gap. (A real deflection tube uses a kV supply and
  can implode, or break inwards: teacher demonstration only.)}
\results{(1)~A graph of $y$ against $V$ is a straight line through the origin. Half the
  speed gives four times the deflection. At the same entry speed, a proton is
  deflected the other way, and about 1840 times less (its mass is about 1840 times
  bigger). $\frac{y}{x^2}$ is the same at every $x$. (2)~$v^2$ against $V$ is a
  straight line through the origin. Its gradient agrees with
  $\frac{2e}{m_\mathrm{e}} = \qty{3.52e11}{C.kg^{-1}}$. The gap makes no difference.}
\conclusion{(1)~So $y \propto V$ and $y \propto x^2$: a parabola, like a projectile,
  because the force is constant. (2)~$v^2 = \frac{2q}{m}V$ is confirmed; $W = qV$
  does not depend on $d$.}
\significance{The same physics steered the electron beam in cathode-ray
  oscilloscopes (old screens that drew voltage graphs). Electron guns, X-ray tubes and accelerators set particle energy
  with the voltage.}
\end{experiment}

%% ==================================================================
\begin{summarybox}{Remember}
\begin{itemize}
  \item Between plates $E = \frac{V}{d}$, $F = qE$ and $a = \frac{qE}{m}$ are
        constant. $W = qV = qEd$; from rest $qV = \frac{1}{2}mv^2$.
  \item Between the plates: a projectile path (parabola) with $\frac{qE}{m}$ in
        place of $g$. After the plates: a straight line.
  \item $F = qvB\sin\theta$ acts only on moving charges, at \ang{90} to $v$ and
        $B$ (palm rule; reverse it for negative charges).
  \item At \ang{90} to $B$: a circle, from $qvB = \frac{mv^2}{r}$:
        $r = \frac{mv}{qB}$, $T = \frac{2\pi m}{qB}$ (not on the sheet).
  \item An electric field can change the speed; a magnetic field does no work and
        changes only the direction.
\end{itemize}
\end{summarybox}
}{}% Charged Particles in Electric and Magnetic Fields
\IfFileExists{./ch2.tex}{\setcounter{chapter}{1}% ch2.tex -- Module 6 ESSENTIALS edition, Chapter 2: The Motor Effect
% Knowledge points only: no worked examples, no questions, no answers, no exam tips.
% Condensed from the final full chapter src/mod6/ch2.tex (same notation, same facts,
% same diagrams; TikZ copied from it, captions shortened).
% Syllabus: NSW Physics Stage 6 (2017), Module 6, IQ2 (M6.IQ2.01-03); Year 11 field refresher.
% Directions checked with vectors: +x right, +y up the page, +z out of the page;
% force on a wire = I l x B; field of a wire = right-hand grip rule.
\tikzset{
  chTwocurrent/.style={->, line width=1.4pt, color=hscink},
  chTwowire/.style={line width=2.2pt, color=hscink!70},
  chTwoloop/.style={line width=1.1pt, color=hscink, midarrow},
  chTwowireout/.pic={\draw[line width=0.9pt, fill=white] (0,0) circle (0.2);
    \fill (0,0) circle (0.06);},
  chTwowirein/.pic={\draw[line width=0.9pt, fill=white] (0,0) circle (0.2);
    \draw[line width=0.9pt] (-0.13,-0.13) -- (0.13,0.13) (-0.13,0.13) -- (0.13,-0.13);},
  chTwofout/.pic={\fill[hscfield] (0,0) circle (0.055);},
  chTwofin/.pic={\draw[hscfield, line width=0.75pt] (-0.085,-0.085) -- (0.085,0.085)
    (-0.085,0.085) -- (0.085,-0.085);},
  chTwoFout/.pic={\draw[hscforce, line width=1.1pt, fill=white] (0,0) circle (0.2);
    \fill[hscforce] (0,0) circle (0.06);},
  chTwoarc/.style={line width=0.8pt, color=hscfield, postaction={decorate},
    decoration={markings, mark=at position 0.15 with {\arrow{Stealth[length=2.2mm, width=1.7mm]}}}},
  chTwoFin/.pic={\draw[hscforce, line width=1.1pt, fill=white] (0,0) circle (0.2);
    \draw[hscforce, line width=1.1pt] (-0.13,-0.13) -- (0.13,0.13) (-0.13,0.13) -- (0.13,-0.13);},
}

\chapter{The Motor Effect}\label{ch2:chap}
\covers{M6.IQ2.01, M6.IQ2.02, M6.IQ2.03}

\begin{summarybox}{In this chapter}
\begin{itemize}
  \item Currents make magnetic fields: around a straight wire, a loop and a solenoid.
  \item The motor effect: $F = lIB\sin\theta$ on a wire in an external magnetic field.
  \item Two parallel wires attract or repel. This force once defined the ampere.
  \item Two key investigations: the current balance and two parallel wires.
\end{itemize}
\end{summarybox}

%% ==================================================================
\section{Magnetic fields made by currents (Year 11 refresher)}\label{ch2:sec:fields}

\begin{bigidea}
Every current makes a magnetic field. Around a straight wire the field lines are circles. A loop
or a solenoid makes a field like a bar magnet's.
\end{bigidea}

\begin{keywords}
  \kw{Magnetic field ($B$)}{where a magnet or moving charge feels a force. Unit: tesla (T).}
  \kw{Conventional current}{the way positive charge would flow: out of the $+$ terminal, round
    the circuit, back to $-$.}
  \kw{Solenoid}{a long coil of wire with many turns.}
  \kw{Uniform field}{the same strength and direction everywhere.}
  \kw{Permeability constant ($\mu_0$)}{links a current to the field it makes in a vacuum.}
\end{keywords}

\begin{itemize}
  \item \term{Field lines} (\cref{ch2:fig:fields}) show the way a compass N end points; closer
        lines = stronger field. Dot ($\bullet$) = out of the page; cross ($\times$) = into the page.
\end{itemize}

\begin{figure}
  \centering
  \begin{tikzpicture}[x=1cm, y=1cm]
    \tikzset{chTwolab/.style={diagramlabel, font=\footnotesize, align=center}}
    % (a) straight wire, end-on
    \begin{scope}
      \foreach \r in {0.35,0.75,1.35} {\draw[field] (\r,0) arc (0:360:\r);}
      \pic at (0,0) {chTwowireout};
      \node[diagramlabel, text=hscfield] at (1.25,1.2) {$B$};
      \node[chTwolab] at (0,-1.75) {(a) wire, end-on:\\ $I$ out of the page};
    \end{scope}
    % (b) straight wire, side view
    \begin{scope}[xshift=3.2cm]
      \draw[chTwowire] (0,-1.3) -- (0,1.3);
      \draw[chTwocurrent] (0,-0.55) -- (0,0.55);
      \node[diagramlabel] at (0,1.52) {$I$};
      \foreach \y in {-1.1,-0.55,0,0.55,1.1} {\pic at (-0.32,\y) {chTwofout}; \pic at (0.32,\y) {chTwofin};}
      \foreach \y in {-0.85,-0.28,0.28,0.85} {\pic at (-0.7,\y) {chTwofout}; \pic at (0.7,\y) {chTwofin};}
      \foreach \y in {-0.6,0.6} {\pic at (-1.08,\y) {chTwofout}; \pic at (1.08,\y) {chTwofin};}
      \node[chTwolab] at (0,-1.75) {(b) wire, side view:\\ $I$ up the page};
    \end{scope}
    % (c) loops, face-on: anticlockwise = N face, clockwise = S face
    \begin{scope}[xshift=6.35cm]
      \draw[chTwoloop] (-0.15,0.1) ++(270:0.58) arc (270:630:0.58);
      \node[diagramlabel, font=\footnotesize, anchor=south] at (-0.15,0.84) {$I$};
      \foreach \p in {(-0.15,0.1),(-0.38,0.1),(0.08,0.1),(-0.15,0.33),(-0.15,-0.13)} {\pic at \p {chTwofout};}
      \node[diagramlabel, font=\footnotesize\bfseries, text=hscforce] at (-0.15,-0.78) {N face};
      \begin{scope}[xshift=1.45cm]
        \draw[chTwoloop] (-0.15,0.1) ++(270:0.58) arc (270:-90:0.58);
        \node[diagramlabel, font=\footnotesize, anchor=south] at (-0.15,0.84) {$I$};
        \foreach \p in {(-0.15,0.1),(-0.38,0.1),(0.08,0.1),(-0.15,0.33),(-0.15,-0.13)} {\pic at \p {chTwofin};}
        \node[diagramlabel, font=\footnotesize\bfseries, text=hscvel] at (-0.15,-0.78) {S face};
      \end{scope}
      \node[chTwolab] at (0.58,-1.75) {(c) loops, face-on:\\ N face and S face};
    \end{scope}
    % (d) solenoid cut along its axis
    \begin{scope}[xshift=10.95cm]
      \fill[hscbrand!6] (-1.5,-0.42) rectangle (1.5,0.42);
      \foreach \x in {-1.25,-0.75,...,1.3} {\pic at (\x,0.42) {chTwowireout}; \pic at (\x,-0.42) {chTwowirein};}
      \foreach \s in {1,-1} {
        \draw[hscfield, line width=0.8pt, postaction={decorate},
          decoration={markings, mark=at position 0.16 with {\arrow{Stealth[length=2.2mm, width=1.7mm]}},
          mark=at position 0.64 with {\arrow{Stealth[length=2.2mm, width=1.7mm]}}}]
          (-1.6,0.17*\s) -- (1.6,0.17*\s) .. controls (2.5,0.25*\s) and (2.45,1.3*\s) .. (0,1.3*\s)
          .. controls (-2.45,1.3*\s) and (-2.5,0.25*\s) .. cycle;
      }
      \node[font=\small\bfseries, text=hscforce] at (1.78,0) {N};
      \node[font=\small\bfseries, text=hscvel] at (-1.78,0) {S};
      \node[chTwolab] at (0,-1.75) {(d) solenoid, cut in half:\\ $I$ out on top, in below};
    \end{scope}
  \end{tikzpicture}
  \caption{Fields made by currents. Inside a solenoid the field is uniform, from S to
    N.}\label{ch2:fig:fields}
\end{figure}

\begin{sidebyside}[0.34]
  \centering
  \begin{tikzpicture}[x=1cm, y=1cm]
    \tikzset{hand/.style={draw=hscink, line width=0.8pt, fill=hscbrand!7}}
    % back half of a field line (behind the wire), then the wire
    \draw[hscfield, line width=0.9pt, postaction={decorate}, decoration={markings,
      mark=at position 0.3 with {\arrow{Stealth[length=2.2mm, width=1.7mm]}}}]
      (0.95,-1.05) arc (0:180:0.95 and 0.3);
    \draw[chTwowire] (0,-1.7) -- (0,1.75);
    % front half of the field line, with its arrows (away from the wire)
    \draw[hscfield, line width=0.9pt, postaction={decorate}, decoration={markings,
      mark=at position 0.3 with {\arrow{Stealth[length=2.2mm, width=1.7mm]}},
      mark=at position 0.75 with {\arrow{Stealth[length=2.2mm, width=1.7mm]}}}]
      (-0.95,-1.05) arc (180:360:0.95 and 0.3);
    % fist (curled fingers) and thumb
    \draw[hand, rounded corners=8pt] (-0.6,-0.55) rectangle (0.6,0.5);
    \foreach \y in {-0.29,-0.03,0.23} {\draw[hscink, line width=0.6pt] (-0.6,\y) .. controls (0,\y-0.1) .. (0.6,\y);}
    \draw[hand, rounded corners=5pt] (-0.17,0.32) rectangle (0.17,1.25);
    \draw[chTwocurrent] (0,1.3) -- (0,1.95) node[above, inner sep=2pt] {$I$};
    \node[diagramlabel, anchor=east, align=right] at (-0.3,0.9) {right\\ hand};
    \node[diagramlabel, align=left, anchor=west] at (0.32,0.95) {thumb:\\ current $I$};
    \node[diagramlabel, align=left, anchor=west, text=hscfield] at (1.0,-1.05) {curled fingers:\\ field $B$};
  \end{tikzpicture}
  \captionof{figure}{The right-hand grip rule.}\label{ch2:fig:grip}
\nextside
  \begin{itemize}
    \item \textbf{Straight wire:} use the \term{right-hand grip rule} (\cref{ch2:fig:grip}).
          Right thumb along the conventional current; the curled fingers show the field.
    \item \textbf{Current loop} (one turn): curl your right fingers with the current; the thumb
          points out of the N face. Face-on: anticlockwise = N, clockwise = S.
    \item \textbf{Solenoid:} the fields of the loops add. Inside: strong, uniform, from S to N.
          Outside: like a bar magnet. The thumb points to the N end.
  \end{itemize}
\end{sidebyside}


\begin{formulacard}[unbreakable]{Fields of a long straight wire and a solenoid}
  \cardformula{B = \frac{\mu_0 I}{2\pi r} \qquad\qquad B = \frac{\mu_0 NI}{L}}
  \begin{symbols}
    \sym{$B$}{magnetic field strength}{\si{T}}
    \sym{$\mu_0$}{magnetic permeability constant, $4\pi \times 10^{-7}$}{\si{N.A^{-2}}}
    \sym{$I$}{current}{\si{A}}
    \sym{$r$}{wire: shortest distance from the centre of the wire}{\si{m}}
    \sym{$N$, $L$}{solenoid: number of turns; length of the solenoid}{no unit; \si{m}}
  \end{symbols}
  \whentouse{a long straight wire (left); inside a long air-core solenoid (right).}
  \trap{wire: $B \propto 1/r$, not $1/r^2$. Solenoid: $L$ is the coil's length, not the wire's.}
\end{formulacard}

%% ==================================================================
\section{The motor effect: force on a wire}\label{ch2:sec:motor}
\covers{M6.IQ2.01}

\begin{bigidea}
A wire carrying a current in an external magnetic field feels a force. The force is largest when
the wire is at \ang{90} to the field. There is no force when the wire is parallel to the field.
\end{bigidea}

\begin{keywords}
  \kw{Motor effect}{the force on a current-carrying wire in an \emph{external} magnetic field.}
  \kw{External field}{a field from something else (such as a magnet), not from the wire.}
\end{keywords}

\begin{itemize}
  \item \textbf{Why:} a current is a stream of moving charges. A moving charge in a magnetic field
        feels a force ($F = qv_\perp B$). The charges cannot leave the wire, so the wire is pushed.
  \item $F$ grows with $l$ (the length in the field), $I$ and $B$. Only the part of the current at
        \ang{90} to $B$ counts: $I_\perp = I\sin\theta$, with $\theta$ between the wire and $B$
        (\cref{ch2:fig:cases}).
\end{itemize}

\begin{figure}
  \centering
  \begin{tikzpicture}[x=1cm, y=1cm]
    \foreach \dx in {0, 4.7, 9.4} {
      \begin{scope}[xshift=\dx cm]
        \foreach \y in {-1.0,-0.33,0.33,1.0} {\draw[field, hscfield!70] (-1.7,\y) -- (1.7,\y);}
        \node[diagramlabel, text=hscfield] at (1.9,1.0) {$B$};
      \end{scope}
    }
    % (a) perpendicular
    \draw[chTwowire] (0,-1.3) -- (0,1.3);
    \draw[chTwocurrent] (0,-0.25) -- (0,0.7) node[left, inner sep=3pt] {$I$};
    \pic at (0.55,-0.67) {chTwoFin};
    \node[diagramlabel, text=hscforce, anchor=west] at (0.78,-0.67) {$F$};
    \node[diagramlabel, align=center] at (0,-2.05) {(a) $\theta = \ang{90}$:\\ largest force, $F = lIB$};
    % (b) parallel
    \begin{scope}[xshift=4.7cm]
      \draw[chTwowire] (-1.4,0) -- (1.4,0);
      \draw[chTwocurrent] (-0.6,0) -- (0.5,0);
      % label at the end of the wire, between two field lines (it does not touch them)
      \node[anchor=west, inner sep=2pt] at (1.4,0) {$I$};
      \node[diagramlabel, align=center] at (0,-2.05) {(b) $\theta = \ang{0}$:\\ no force};
    \end{scope}
    % (c) at an angle
    \begin{scope}[xshift=9.4cm]
      \coordinate (A) at (-0.95,-0.8); \coordinate (B) at (0.95,0.8);
      \draw[chTwowire] ($(A)!-0.12!(B)$) -- ($(A)!1.12!(B)$);
      \draw[chTwocurrent] (A) -- (B) node[pos=0.62, below right, inner sep=1pt, fill=white] {$I$};
      \draw[->, line width=1pt, hscacc] (A) -- (0.95,-0.8) node[pos=0.5, below=9pt, inner sep=1pt, text=hscacc] {$I_\parallel = I\cos\theta$};
      \draw[->, line width=1pt, hscacc] (0.95,-0.8) -- (B) node[pos=0.5, right, inner sep=2pt, text=hscacc, fill=white] {$I_\perp = I\sin\theta$};
      \draw[hscink, line width=0.6pt] ($(A)+(0.55,0)$) arc (0:40.1:0.55);
      \node[diagramlabel, fill=white, inner sep=1pt] at ($(A)+(0.78,0.26)$) {$\theta$};
      \node[diagramlabel, align=center] at (0,-2.05) {(c) angle $\theta$: only $I_\perp$\\ counts, $F = lIB\sin\theta$};
    \end{scope}
  \end{tikzpicture}
  \caption{A wire in a uniform field that points to the right. In (a) the force is into the page.}\label{ch2:fig:cases}
\end{figure}

\begin{itemize}
  \item \textbf{Largest force:} wire at \ang{90} to $B$ ($\sin\ang{90} = 1$): $F = lIB$. In
        between, $F$ follows $\sin\theta$.
  \item \textbf{No force:} wire parallel to $B$ ($\theta = \ang{0}$ or \ang{180}), no current, or
        wire outside the field.
  \item \textbf{Special case:} $B$ into or out of the page: every wire in the page is at \ang{90}
        to $B$.
  \item \textbf{Direction:} $F$ is at \ang{90} to both $I$ and $B$. Use the
        \term{right-hand palm rule} (\cref{ch2:fig:palm}).
\end{itemize}

\begin{sidebyside}[0.42]
  \centering
  \begin{tikzpicture}[x=0.85cm, y=0.85cm]
    \tikzset{hand/.style={draw=hscink, line width=0.8pt, fill=hscbrand!7}}
    % thumb, fingers, wrist, then palm on top
    \draw[hand, rounded corners=6pt] (1.45,0.32) rectangle (2.95,0.82);
    \draw[hand, rounded corners=5pt] (0.05,1.6) rectangle (0.43,2.8);
    \draw[hand, rounded corners=5pt] (0.47,1.6) rectangle (0.86,3.2);
    \draw[hand, rounded corners=5pt] (0.90,1.6) rectangle (1.29,3.35);
    \draw[hand, rounded corners=5pt] (1.33,1.6) rectangle (1.72,3.12);
    \draw[hand] (0.3,-0.75) rectangle (1.45,0.2);
    \draw[hand, rounded corners=9pt] (0,0) rectangle (1.77,1.9);
    \pic at (0.88,0.95) {chTwoFout};
    % arrows and labels
    \draw[->, line width=1.3pt, hscfield] (-0.45,1.4) -- (-0.45,3.3);
    \node[diagramlabel, text=hscfield, align=right, anchor=east] at (-0.55,2.4) {fingers:\\ field $B$};
    \draw[chTwocurrent] (3.1,0.57) -- (4.1,0.57);
    \node[diagramlabel, align=center, anchor=north] at (3.6,0.4) {thumb:\\ current $I$};
    \node[diagramlabel, text=hscforce, align=left, anchor=west] at (2.0,1.75) {palm pushes:\\ force $F$ out\\ of the page};
    \draw[hscforce, line width=0.5pt] (1.97,1.6) -- (1.1,1.05);
  \end{tikzpicture}
  \captionof{figure}{The right-hand palm rule.}\label{ch2:fig:palm}
\nextside
  \begin{enumerate}
    \item Hold your right hand flat.
    \item \textbf{Fingers} along the field $B$ (N to S).
    \item \textbf{Thumb} along the conventional current $I$ (not the electron flow).
    \item The \textbf{palm} faces the force $F$.
  \end{enumerate}
  Reverse $I$ \emph{or} $B$: $F$ reverses. Reverse both: no change.
\end{sidebyside}

\begin{formulacard}[unbreakable]{Force on a current-carrying wire}
  \cardformula{F = lI_\perp B = lIB\sin\theta}
  \begin{symbols}
    \sym{$F$}{size of the force on the wire}{\si{N}}
    \sym{$l$}{length of wire inside the field}{\si{m}}
    \sym{$I$}{current in the wire; $I_\perp = I\sin\theta$ is the part at \ang{90} to $B$}{\si{A}}
    \sym{$B$}{magnetic field strength of the external field}{\si{T} ($= \si{N.A^{-1}.m^{-1}}$)}
    \sym{$\theta$}{angle between the wire (the current) and the field}{degree (\si{\degree})}
  \end{symbols}
  \whentouse{a straight wire with a current in a uniform external magnetic field.}
  \trap{for an angle measured from a line at \ang{90} to $B$, $\theta = \ang{90} - (\text{that angle})$.}
\end{formulacard}

%% ==================================================================
\section{Forces between parallel wires}\label{ch2:sec:wires}
\covers{M6.IQ2.03}

\begin{bigidea}
Each of two parallel wires sits in the field of the other, so each feels a motor-effect force.
Same-direction currents attract; opposite currents repel.
\end{bigidea}

\begin{keywords}
  \kw{Force per unit length ($F/l$)}{the force on each metre of wire. Unit: \si{N.m^{-1}}.}
  \kw{Newton's third law}{if wire 1 pushes wire 2, wire 2 pushes back equally hard.}
\end{keywords}

\begin{itemize}
  \item Direction (\cref{ch2:fig:parallel}): (1) grip rule: field $B_1$ of wire 1 at wire 2;
        (2) palm rule: force on wire 2; (3) Newton's third law: equal and opposite force on wire 1.
  \item \textbf{Same direction: attract} (pull together). \textbf{Opposite: repel} (push apart).
        This is the reverse of electric charges, where like charges repel.
  \item \textbf{Why:} wire 2 is at \ang{90} to the field of wire 1, $B_1 = \mu_0 I_1/(2\pi r)$,
        so $F = lI_2B_1$. $F$ grows with $I_1$ and $I_2$, and falls as the gap $r$ grows.
\end{itemize}

\begin{figure}
  \centering
  \begin{tikzpicture}[x=1cm, y=1cm]
    \foreach \dx in {0, 7.2} {
      \begin{scope}[xshift=\dx cm]
        \draw[guide] (0,0) -- (3,0);
        \draw[chTwoarc] (-24:3) arc (-24:24:3);
        \draw[->, line width=1.2pt, hscfield] (3,0.2) -- (3,1.0) node[right, inner sep=2pt] {$B_1$};
      \end{scope}
    }
    % (a) same direction
    \draw[chTwoarc] (156:3) ++(3,0) arc (156:204:3);
    \draw[->, line width=1.2pt, hscfield] (0,-0.2) -- (0,-1.0) node[left, inner sep=2pt] {$B_2$};
    \pic at (0,0) {chTwowireout}; \pic at (3,0) {chTwowireout};
    \draw[force] (0.25,0) -- (1.2,0) node[midway, above, inner sep=2pt, text=hscforce] {$F$};
    \draw[force] (2.75,0) -- (1.8,0) node[midway, above, inner sep=2pt, text=hscforce] {$F$};
    \node[diagramlabel, align=right, anchor=east] at (-0.3,0.4) {wire 1\\ \qty{2.0}{A}};
    \node[diagramlabel, align=left, anchor=west] at (3.3,-0.4) {wire 2\\ \qty{6.0}{A}};
    \node[diagramlabel] at (1.5,-1.75) {(a) same direction: \textbf{attract}};
    % (b) opposite directions
    \begin{scope}[xshift=7.2cm]
      \draw[chTwoarc] (204:3) ++(3,0) arc (204:156:3);
      \draw[->, line width=1.2pt, hscfield] (0,0.2) -- (0,1.0) node[left, inner sep=2pt] {$B_2$};
      \pic at (0,0) {chTwowireout}; \pic at (3,0) {chTwowirein};
      \draw[force] (-0.25,0) -- (-1.25,0) node[above, inner sep=2pt, text=hscforce] {$F$};
      \draw[force] (3.25,0) -- (4.25,0) node[above, inner sep=2pt, text=hscforce] {$F$};
      \node[diagramlabel, anchor=north east] at (-0.2,-0.25) {wire 1};
      \node[diagramlabel, anchor=north west] at (3.2,-0.25) {wire 2};
      \node[diagramlabel] at (1.5,-1.75) {(b) opposite directions: \textbf{repel}};
    \end{scope}
  \end{tikzpicture}
  \caption{Two wires seen end-on. In (a) the currents differ, but the two forces are the same
    size.}\label{ch2:fig:parallel}
\end{figure}

\begin{formulacard}[unbreakable]{Force between two parallel wires}
  \cardformula{\frac{F}{l} = \frac{\mu_0}{2\pi}\frac{I_1 I_2}{r}}
  \begin{symbols}
    \sym{$F$}{force on \emph{each} wire}{\si{N}}
    \sym{$l$}{length over which the wires run side by side}{\si{m}}
    \sym{$\mu_0$}{magnetic permeability constant; $\mu_0/2\pi = 2 \times 10^{-7}$}{\si{N.A^{-2}}}
    \sym{$I_1$, $I_2$}{the currents in the two wires}{\si{A}}
    \sym{$r$}{distance between the wires (centre to centre)}{\si{m}}
  \end{symbols}
  \whentouse{two long, straight, parallel wires (much longer than the gap $r$).}
  \trap{$F/l$ is in \si{N.m^{-1}}: multiply by $l$ to get the force. It is $r$, not $r^2$.}
\end{formulacard}

\subsection{Newton's third law and the SI ampere}

\begin{itemize}
  \item Both wires feel the \emph{same size} of force, even when the currents differ ($I_1 I_2$ is
        the same for both). They form an \term{action--reaction pair}: equal and opposite forces
        that two objects put on each other.
  \item The \term{ampere} (A) is the \term{SI} unit (the international unit) of current
        (\qty{1}{A} = \qty{1}{C} per second). From 1948 to 2019 it was defined by this force
        (since 2019, by the size of the electron's charge).
        \textbf{Learn the wording:} \emph{one ampere is the constant current that, flowing in
        each of two straight, parallel wires of infinite length and negligible circular
        cross-section} (very thin)\emph{, placed \qty{1}{m} apart in a vacuum} (empty
        space)\emph{, produces \textbf{between} these wires a force of $2 \times 10^{-7}$
        newton per metre of length.}
  \item \textbf{Link to Newton's third law:} the definition names one force \emph{between} the
        wires, not which wire. Newton's third law makes the two forces equal and opposite, so one
        value is true for both wires.
\end{itemize}

%% ==================================================================
\section{Key experiments}\label{ch2:sec:experiments}
\covers{M6.IQ2.01, M6.IQ2.02}

Both investigations use an \term{electronic balance} (a digital scale, \cref{ch2:fig:apparatus}).
\begin{figure}
  \centering
  \begin{tikzpicture}[x=1cm, y=1cm]
    % ---- (a) current balance
    \draw[object, rounded corners=3pt] (0,0) rectangle (5.6,0.7);
    \draw[hscink, fill=white, line width=0.6pt] (0.25,0.15) rectangle (1.55,0.55);
    \node[font=\scriptsize] at (0.9,0.35) {\num{0.00} g};
    \draw[object] (1.4,0.7) rectangle (4.4,0.85);
    \node[diagramlabel, anchor=west] at (2.0,0.35) {electronic balance};
    \draw[draw=hscink, fill=hscmuted!25, line width=0.8pt] (1.7,0.85) -- (4.1,0.85) -- (4.1,2.6)
      -- (3.6,2.6) -- (3.6,1.25) -- (2.2,1.25) -- (2.2,2.6) -- (1.7,2.6) -- cycle;
    \draw[draw=hscink, fill=hscforce!20, line width=0.8pt] (2.2,1.5) rectangle (2.52,2.4);
    \draw[draw=hscink, fill=hscvel!20, line width=0.8pt] (3.28,1.5) rectangle (3.6,2.4);
    \node[font=\small\bfseries] at (2.36,1.95) {N};
    \node[font=\small\bfseries] at (3.44,1.95) {S};
    \foreach \y in {1.63,2.27} {\draw[field] (2.54,\y) -- (3.26,\y);}
    \pic at (2.9,1.95) {chTwowireout};
    % clamp holding the wire (the wire does not touch the magnet)
    \draw[line width=2pt, hscmuted] (1.0,3.05) -- (2.9,3.05) -- (2.9,2.15);
    \node[diagramlabel, anchor=east] at (0.95,3.05) {clamp};
    \draw[force] (3.1,2.2) -- (3.1,2.85) node[right, inner sep=2pt, text=hscforce] {$F$ on wire};
    \draw[force] (4.35,1.9) -- (4.35,1.2);
    \node[diagramlabel, text=hscforce, anchor=west, align=left] at (4.4,1.55) {$F$ on\\ magnet};
    \node[diagramlabel] at (2.8,-0.45) {(a) current balance};
    % ---- (b) two parallel rods on a balance
    \begin{scope}[xshift=7.6cm]
      \tikzset{chTwosmall/.style={diagramlabel, font=\footnotesize}}
      \draw[object, rounded corners=3pt] (0,0) rectangle (4.3,0.6);
      \draw[hscink, fill=white, line width=0.6pt] (0.2,0.12) rectangle (1.25,0.48);
      \node[font=\scriptsize] at (0.725,0.3) {\num{0.000} g};
      \node[chTwosmall, anchor=west] at (1.35,0.3) {electronic balance};
      \draw[object] (0.4,0.6) rectangle (3.9,0.72);
      \draw[draw=hscink, fill=hscmuted!30] (0.7,0.72) rectangle (1.0,1.18);
      \draw[draw=hscink, fill=hscmuted!30] (3.3,0.72) rectangle (3.6,1.18);
      \draw[line width=3pt, hscmuted!80!black] (0.3,1.25) -- (4.0,1.25);
      \draw[line width=3pt, hscmuted!80!black] (0.3,2.05) -- (4.0,2.05);
      \draw[line width=0.8pt] (4.0,2.05) .. controls (4.55,2.05) and (4.55,1.25) .. (4.0,1.25);
      \node[chTwosmall, anchor=west, align=left] at (4.45,1.65) {flexible\\ link};
      \draw[line width=2pt, hscmuted] (1.8,2.12) -- (1.8,2.95);
      \node[chTwosmall, anchor=west] at (1.9,2.78) {insulated clamp};
      \draw[line width=0.8pt] (0.3,2.05) -- (-0.1,2.05) (0.3,1.25) -- (-0.1,1.25);
      \fill (-0.1,2.05) circle (1.3pt); \fill (-0.1,1.25) circle (1.3pt);
      \node[chTwosmall, anchor=east, align=right] at (-0.15,1.65) {to\\ supply};
      \draw[chTwocurrent] (0.45,2.27) -- (1.2,2.27) node[right, inner sep=2pt] {$I$};
      \draw[chTwocurrent] (3.1,1.0) -- (2.4,1.0) node[left, inner sep=2pt] {$I$};
      \draw[<->, hscmuted, line width=0.5pt] (0.3,3.2) -- (4.0,3.2);
      \node[chTwosmall, fill=white, inner sep=1pt] at (2.15,3.2) {$l = \qty{0.50}{m}$};
      \draw[<->, hscmuted, line width=0.5pt] (2.75,1.33) -- (2.75,1.97);
      \node[chTwosmall, anchor=west] at (2.8,1.65) {$r$};
      \draw[force] (1.5,1.15) -- (1.5,0.78);
      \node[chTwosmall, text=hscforce, anchor=west] at (1.55,0.93) {$F$};
      \node[diagramlabel] at (2.15,-0.45) {(b) two parallel rods};
    \end{scope}
  \end{tikzpicture}
  \caption{(a) Wire end-on, current out of the page. (b) Rods in series: opposite
    currents.}\label{ch2:fig:apparatus}
\end{figure}

\begin{experiment}{The force on a wire (current balance)}
\aim{To find how the force on a wire depends on the current and the angle.}
\method{\begin{itemize}
  \item A magnet sits on the balance (\cref{ch2:fig:apparatus}(a)). A clamped wire passes through
        the gap at \ang{90} to $B$, without touching the magnet. A DC supply, ammeter,
        \term{rheostat} (a variable resistor) and switch are in series with it.
  \item \term{Tare} the balance (set it to zero). Change $I$ (keep $l$, $B$, $\theta$ the same).
        Reverse $I$, then $B$. Then change $\theta$.
  \item The wire is pushed up (palm rule), so the magnet is pushed down (Newton's third law). The
        reading rises by $\Delta m$, and $F = \Delta m\, g$ ($\Delta m$ in kg).
\end{itemize}}
\results{$F$ against $I$ is a straight line through the origin: $F \propto I$, gradient $lB$.
  $F$ is largest at \ang{90} and zero when the wire is parallel to $B$ (it follows $\sin\theta$).
  Reversing $I$ \emph{or} $B$ reverses $F$; reversing both does not.}
\conclusion{This agrees with $F = lIB\sin\theta$ and the palm rule. Also, $B = \text{gradient} \div l$.}
\significance{Electric motors, loudspeakers and analogue meters all use this effect.}
\end{experiment}

\begin{experiment}{Two parallel wires (required investigation)}
\aim{To show that parallel currents attract or repel, and to test $F/l = (\mu_0/2\pi)(I_1 I_2/r)$.}
\method{\begin{itemize}
  \item The lower rod rests on insulating blocks on the balance; the upper rod is clamped above it
        (\cref{ch2:fig:apparatus}(b)). A flexible link joins them, so the currents are opposite and
        $I_1 I_2 = I^2$. Tare the balance (set it to zero).
  \item A low-voltage, high-current DC supply, ammeter, rheostat and switch are in series with the
        rods; use short bursts of current. $r$ (centre to centre) = the gap, measured with
        vernier calipers (a tool for small lengths), + one rod diameter.
  \item Change $I$; keep $l$ and $r$ the same. Plot $F = \Delta m\, g$ against $I^2$. Change $r$
        with spacers to test $F \propto 1/r$.
  \item Without the link, both rods can carry current the same way.
\end{itemize}}
\results{Opposite currents: the reading rises (repel). Same direction: it falls below zero
  (attract). Hanging foil strips show the same. $F$ against $I^2$ is a straight line through the
  origin: $F \propto I_1 I_2$, gradient $\mu_0 l/(2\pi r)$. This gives $\mu_0$ close to
  $4\pi \times 10^{-7}$~\si{N.A^{-2}}.}
\conclusion{Like currents attract, opposite currents repel, and $F \propto I_1 I_2$, as predicted.}
\significance{Ampère found this force in 1820. It defined the ampere (1948--2019).}
\end{experiment}

%% ==================================================================
\begin{summarybox}{Remember}
\begin{itemize}
  \item Wire: circles, $B = \mu_0 I/(2\pi r)$. Solenoid: uniform inside, $B = \mu_0 NI/L$.
  \item Grip rule, straight wire: thumb = conventional $I$, curled fingers = $B$. Coil: curled
        fingers = current, thumb points to the N end.
  \item Motor effect: $F = lIB\sin\theta$ in an \emph{external} field: largest at \ang{90}, zero
        when parallel.
  \item Palm rule: fingers $B$, thumb conventional $I$, palm $F$ (at \ang{90} to both).
  \item Parallel wires: same direction attract, opposite repel; $F/l = (\mu_0/2\pi)(I_1 I_2/r)$.
  \item Equal and opposite forces (Newton's third law): this force defined the ampere.
\end{itemize}
\end{summarybox}
}{}% The Motor Effect
\IfFileExists{./ch3.tex}{\setcounter{chapter}{2}% ch3.tex -- Module 6 ESSENTIALS, Chapter 3: Electromagnetic induction (flux, Faraday, Lenz)
% Condensed from src/mod6/ch3.tex (final full book). Knowledge points only: no questions,
% no worked examples, no exam tips. Same notation and facts as the full chapter.
% Figure files: figures/ch3_falling.pdf (+ .py), copied from the full book. All other figures
% are TikZ copied (some shrunk) from the full chapter.
\chapter{Electromagnetic Induction: Flux, Faraday and Lenz}\label{ch3:chap}
\covers{M6.IQ3.01, M6.IQ3.02}
% Local TikZ helpers (chapter prefix): a small open loop with one arrowhead.
% \chThreeACW{x}{y}{r} = anticlockwise loop, \chThreeCW{x}{y}{r} = clockwise loop.
\newcommand\chThreeACW[3]{%
  \draw[line width=1.2pt, hscink] ($(#1,#2)+(110:#3)$) arc (110:400:#3);
  \draw[line width=1.2pt, hscink, -{Latex[length=2.4mm, width=2.0mm]}]
    ($(#1,#2)+(40:#3)+(310:0.05)$) -- ($(#1,#2)+(40:#3)+(130:0.19)$);}
\newcommand\chThreeCW[3]{%
  \draw[line width=1.2pt, hscink] ($(#1,#2)+(70:#3)$) arc (70:-220:#3);
  \draw[line width=1.2pt, hscink, -{Latex[length=2.4mm, width=2.0mm]}]
    ($(#1,#2)+(140:#3)+(230:0.05)$) -- ($(#1,#2)+(140:#3)+(50:0.19)$);}

\begin{summarybox}{In this chapter}
\begin{itemize}
  \item Magnetic flux, $\Phi = B_\parallel A = BA\cos\theta$, and how it can change.
  \item Faraday's law: a changing flux induces an emf.
  \item Lenz's law: the induced current opposes the change (conservation of energy).
  \item Induction in moving rods, solenoids, metal plates and two nearby coils.
\end{itemize}
\end{summarybox}

%% ===================================================================
\section{Magnetic flux}\label{ch3:sec:flux}
\covers{M6.IQ3.01}

\begin{bigidea}
Magnetic flux tells you how much magnetic field passes through a loop. It
depends on the field strength, the area and the angle.
\end{bigidea}

\begin{keywords}
  \kw{Magnetic flux ($\Phi$)}{how much magnetic field passes through an area.
    Unit: weber (Wb).}
  \kw{Normal}{an imaginary arrow at \ang{90} to the flat surface of a loop.}
\end{keywords}

\begin{figure}
  \centering
  \begin{tikzpicture}[x=1cm, y=1cm]
    % ---------- (a) theta = 0
    \begin{scope}
      \foreach \y in {0.2, 0.75, 1.8, 2.35} \draw[field] (0,\y) -- (3.6,\y);
      \node[diagramlabel, text=hscfield, right] at (3.6,2.35) {$B$};
      \coordinate (C) at (1.5,1.275);
      \draw[line width=2.4pt, hscbrand] ($(C)+(90:0.7)$) -- ($(C)+(270:0.7)$);
      \draw[->, line width=1.1pt, hscink] (C) -- ++(0:1.0)
        node[pos=1, right, font=\footnotesize, inner sep=2pt] {normal};
      \node[align=center, font=\small] at (1.8,-0.5)
        {(a) $\theta = \ang{0}$\\ $\Phi = BA$ (largest)};
    \end{scope}
    % ---------- (b) theta = 50 deg
    \begin{scope}[xshift=4.75cm]
      \foreach \y in {0.2, 0.75, 1.8, 2.35} \draw[field] (0,\y) -- (3.6,\y);
      \node[diagramlabel, text=hscfield, right] at (3.6,2.35) {$B$};
      \coordinate (C) at (1.5,1.275);
      \draw[guide] (C) -- ++(0:1.1);
      \draw[line width=2.4pt, hscbrand] ($(C)+(140:0.7)$) -- ($(C)+(320:0.7)$);
      \draw[->, line width=1.1pt, hscink] (C) -- ++(50:0.9)
        node[anchor=south west, font=\footnotesize, inner sep=1pt] {normal};
      \draw[line width=0.7pt, hscink] ($(C)+(0:0.45)$) arc (0:50:0.45);
      \node[font=\small, inner sep=0.5pt] at ($(C)+(22:0.68)$) {$\theta$};
      \draw[line width=0.8pt, hscmuted] ($(C)+(320:0.55)$) arc (320:360:0.55);
      \node[font=\footnotesize, text=hscmuted, inner sep=1pt, anchor=west]
        at ($(C)+(338:0.8)$) {$\ang{90} - \theta$};
      \node[align=center, font=\small] at (1.8,-0.5)
        {(b) tilted\\ $\Phi = BA\cos\theta$};
    \end{scope}
    % ---------- (c) theta = 90 deg
    \begin{scope}[xshift=9.5cm]
      \foreach \y in {0.2, 0.75, 1.8, 2.35} \draw[field] (0,\y) -- (3.6,\y);
      \node[diagramlabel, text=hscfield, right] at (3.6,2.35) {$B$};
      \coordinate (C) at (1.2,1.275);
      \draw[line width=2.4pt, hscbrand] ($(C)+(180:0.7)$) -- ($(C)+(0:0.7)$);
      \draw[->, line width=1.1pt, hscink] (C) -- ++(90:0.9)
        node[pos=0.89, anchor=east, font=\footnotesize, inner sep=1.5pt, fill=white] {normal};
      \node[align=center, font=\small] at (1.8,-0.5)
        {(c) $\theta = \ang{90}$: plane\\ parallel to $B$, \ $\Phi = 0$};
    \end{scope}
  \end{tikzpicture}
  \caption{A flat loop in a uniform field, seen from the side (thick navy
    line).}\label{ch3:fig:angle}
\end{figure}

\begin{itemize}
  \item The \term{plane} of a loop is the flat surface inside it, like the skin of a
        drum. $\theta$ is measured from the \emph{normal}, not from the plane
        (\cref{ch3:fig:angle}).
  \item Only the part of $B$ along the normal, $B_\parallel = B\cos\theta$, passes
        through. For an angle to the plane, use $\theta = \ang{90} - \text{that angle}$.
  \item $\qty{1}{Wb} = \qty{1}{T.m^2}$. With $B$ perpendicular to the plane,
        $B = \dfrac{\Phi}{A}$, the flux per square metre. So $B$ is also called the
        \term{magnetic flux density}.
  \item Reversing the field (or turning the loop over) takes the flux from $+BA$ to
        $-BA$. So $\lvert\Delta\Phi\rvert = 2BA$, not zero.
\end{itemize}


\begin{formulacard}{Magnetic flux}
  \cardformula{\Phi = B_\parallel A = BA\cos\theta}
  \begin{symbols}
    \sym{$\Phi$}{magnetic flux through the loop}{\si{Wb} ($= \si{T.m^2}$)}
    \sym{$B_\parallel$}{the part of $B$ along the normal of the loop}{\si{T}}
    \sym{$B$}{magnetic field strength (flux density)}{\si{T}}
    \sym{$A$}{area of the loop that is inside the field}{\si{m^2}}
    \sym{$\theta$}{angle between $B$ and the \emph{normal} to the loop}{\si{\degree}
      (SI unit: rad)}
  \end{symbols}
  \whentouse{a flat loop or coil in a uniform field. $\Phi$ is for \emph{one} turn.}
  \trap{plane \emph{parallel} to the field: $\theta = \ang{90}$, so $\Phi = 0$.
    Change \si{cm^2} to \si{m^2} ($\times\,10^{-4}$).}
\end{formulacard}

\subsection{Three ways to change the flux}
\begin{itemize}
  \item \textbf{Change $B$}: move a magnet closer or away; change the current in an
        electromagnet; switch the field on, off, or reverse it.
  \item \textbf{Change $A$}: slide the loop into or out of the field; squash it.
        \textbf{Change $\theta$}: rotate the loop, as in a generator.
  \item A loop sliding \emph{inside} a large uniform field keeps the same flux: no emf.
\end{itemize}

%% ===================================================================
\section{Faraday's law}\label{ch3:sec:faraday}
\covers{M6.IQ3.02}

\begin{bigidea}
Only a \emph{changing} flux induces an emf. A steady flux induces nothing, however
big it is.
\end{bigidea}

\begin{itemize}
  \item \term{Electromagnetic induction}: making (\emph{inducing}) an emf by changing flux.
  \item \term{emf} ($\varepsilon$, electromotive force): the energy given to each
        coulomb of charge (not a force). Unit: volt (V).
  \item A faster magnet, a stronger magnet or more turns give a bigger emf.
  \item \textbf{Faraday's law}: the induced emf is proportional to the rate of
        change of magnetic flux through the circuit.
  \item The $N$ turns add their emfs, like cells in a battery. Closed circuit: an
        \term{induced current} $I = \dfrac{\varepsilon}{R}$ flows. Open circuit: an
        emf, but no current.
  \item On a $\Phi$--$t$ graph, $\lvert\varepsilon\rvert = N \times \text{gradient}$.
        Steep: big emf. Flat: zero. Rising and falling parts give opposite signs.
\end{itemize}

\begin{formulacard}{Faraday's law of induction}
  \cardformula{\varepsilon = -N\frac{\Delta\Phi}{\Delta t}}
  \begin{symbols}
    \sym{$\varepsilon$}{average induced emf}{\si{V}}
    \sym{$N$}{number of turns in the coil}{no unit}
    \sym{$\Delta\Phi$}{change in flux through one turn:
      $\Phi_\mathrm{final} - \Phi_\mathrm{initial}$}{\si{Wb}}
    \sym{$\Delta t$}{time taken for the change}{\si{s}}
  \end{symbols}
  \whentouse{the size of the emf for any change of flux.}
  \trap{$\Delta\Phi$ is the \emph{change} in flux: a huge steady flux gives zero emf.
    $\si{Wb.s^{-1}} = \si{V}$.}
\end{formulacard}

%% ===================================================================
\section{Lenz's law: the direction of the current}\label{ch3:sec:lenz}
\covers{M6.IQ3.02}

\begin{bigidea}
An induced current always fights the change that made it. \term{Lenz's law}
(learn it): \emph{the induced current flows in the direction that makes a magnetic
field opposing the change in flux that produced it.}
\end{bigidea}

\begin{itemize}
  \item Lenz's law is the minus sign in Faraday's law. It follows from conservation
        of energy. Open circuit: the emf acts the way that would drive such a current.
  \item The induced field opposes the \emph{change}, not the field. Flux increasing: it
        points opposite to the original field. Flux decreasing: the \emph{same} way.
  \item \term{Right-hand grip rule}: fingers curl with the current, thumb gives the
        field inside. Anticlockwise (as you see it): an N face. Clockwise: an S face.
  \item Steps for any direction: (1) find which way the field points and whether the
        flux grows or shrinks; (2) Lenz's law gives the induced field; (3) the grip
        rule gives the current. Clockwise or anticlockwise has a meaning only when the
        viewing side is named (for example, seen from the magnet).
  \item \term{Pole method} (\cref{ch3:fig:magnetcoil}): the near face of the coil becomes
        the \emph{same} pole as an approaching pole (it repels), and the
        \emph{opposite} pole to a leaving pole (it attracts).
\end{itemize}

\begin{figure}
  \centering
  \begin{tikzpicture}[x=1cm, y=1cm]
    \foreach \xs in {0, 7.1} {
    \begin{scope}[xshift=\xs cm]
      \fill[hscvel!22] (0,0) rectangle (0.8,0.5);
      \fill[hscforce!22] (0.8,0) rectangle (1.6,0.5);
      \draw[line width=0.8pt, hscink] (0,0) rectangle (1.6,0.5);
      \node[font=\small\bfseries] at (0.4,0.25) {S};
      \node[font=\small\bfseries] at (1.2,0.25) {N};
      \foreach \x in {2.6, 2.85, 3.1, 3.35, 3.6}
        \draw[line width=1pt, hscink] (\x,0.25) ellipse (0.1 and 0.6);
    \end{scope}}
    % ----- (a) approaching
    \draw[velocity] (0.3,0.85) -- (1.3,0.85) node[right, text=hscvel] {$v$};
    \draw[->, hscfield, line width=1.2pt] (3.95,0.25) -- (2.05,0.25);
    \node[font=\footnotesize, text=hscfield, anchor=west] at (3.95,0.25) {induced $B$};
    \node[font=\footnotesize\bfseries, text=hscforce] at (3.1,1.15) {near face: N};
    \node[font=\footnotesize, align=center] at (2.2,-0.95)
      {(a) N pole approaching: the near\\ face becomes N and repels.};
    \chThreeACW{1.2}{-1.9}{0.36}
    \node[font=\footnotesize, align=left, anchor=west] at (1.75,-1.9)
      {seen from the magnet:\\ anticlockwise};
    % ----- (b) moving away
    \begin{scope}[xshift=7.1cm]
      \draw[velocity] (1.3,0.85) -- (0.3,0.85) node[left, text=hscvel] {$v$};
      \draw[->, hscfield, line width=1.2pt] (2.05,0.25) -- (3.95,0.25);
      \node[font=\footnotesize, text=hscfield, anchor=west] at (3.95,0.25) {induced $B$};
      \node[font=\footnotesize\bfseries, text=hscvel] at (3.1,1.15) {near face: S};
      \node[font=\footnotesize, align=center] at (2.2,-0.95)
        {(b) N pole moving away: the near\\ face becomes S and attracts.};
      \chThreeCW{1.2}{-1.9}{0.36}
      \node[font=\footnotesize, align=left, anchor=west] at (1.75,-1.9)
        {seen from the magnet:\\ clockwise};
    \end{scope}
  \end{tikzpicture}
  \caption{The pole method. The coil always opposes the motion of the
    magnet.}\label{ch3:fig:magnetcoil}
\end{figure}

\subsection{Lenz's law and conservation of energy}
\begin{itemize}
  \item If a coil \emph{attracted} an approaching magnet, the magnet would speed up
        and gain energy from nothing. That breaks the \term{law of conservation of
        energy}: energy cannot be created or destroyed.
  \item So the coil pushes back: work done against it $\to$ electrical energy $\to$
        heat. Open circuit: no current, no push back.
\end{itemize}

%% ===================================================================
\section{Motional emf: a straight wire moving in a field}\label{ch3:sec:motional}
\covers{M6.IQ3.02}

\begin{bigidea}
A metal rod moving across a magnetic field has an emf between its ends.
\end{bigidea}

\begin{sidebyside}[0.3]
  \centering
  \begin{tikzpicture}[x=0.8cm, y=0.8cm]
    \foreach \x in {2.1, 2.7} \foreach \y in {0.3, 0.9, 2.1, 2.7}
      \node[text=hscfield, font=\small] at (\x,\y) {$\times$};
    \foreach \y in {0.3, 0.9, 1.5, 2.1, 2.7}
      \node[text=hscfield, font=\small] at (3.3,\y) {$\times$};
    \foreach \x in {0.3, 0.9} \foreach \y in {0.3, 0.9, 2.7}
      \node[text=hscfield, font=\small] at (\x,\y) {$\times$};
    \draw[line width=3.2pt, hscmuted] (1.5,0.15) -- (1.5,2.85);
    \node[draw, circle, inner sep=0.6pt, font=\footnotesize\bfseries, fill=white]
      at (1.5,3.15) {$+$};
    \node[draw, circle, inner sep=0.6pt, font=\footnotesize\bfseries, fill=white]
      at (1.5,-0.15) {$-$};
    \draw[velocity] (1.62,1.2) -- (2.5,1.2);
    \node[text=hscvel, font=\small, fill=white, inner sep=1pt, anchor=west]
      at (2.52,1.2) {$v$};
    \draw[force] (1.2,1.25) -- (1.2,2.35);
    \node[font=\footnotesize, text=hscforce, fill=white, inner sep=1pt, anchor=east]
      at (1.12,1.8) {$F$ on $+q$};
  \end{tikzpicture}
  \captionof{figure}{A rod alone, $B$ into the page.}\label{ch3:fig:rod}
\nextside
  \begin{itemize}
    \item \term{Right-hand palm rule}: fingers along $B$, thumb along $v$ of a
          \emph{positive} charge, palm pushes the way of the force.
    \item Rod moving right (\cref{ch3:fig:rod}): positive charge is pushed up, free
          electrons (they can move in the metal) down. The top end becomes $+$, like a
          battery: \term{motional emf}.
    \item The build-up stops when the electric and magnetic forces balance.
  \end{itemize}
\end{sidebyside}

\begin{sidebyside}[0.43]
  \centering
  \begin{tikzpicture}[x=0.75cm, y=0.75cm]
    \fill[hscvel!12] (3.0,0) rectangle (3.8,3.0);
    \foreach \x in {1.65, 2.4, 4.6, 5.4} \foreach \y in {0.5, 1.5, 2.5}
      \node[text=hscfield, font=\small] at (\x,\y) {$\times$};
    \foreach \y in {0.5, 2.5} \node[text=hscfield, font=\small] at (0.9,\y) {$\times$};
    \foreach \y in {0.5, 2.5} \node[text=hscfield, font=\small] at (3.4,\y) {$\times$};
    \draw[line width=1.5pt, hscink] (0,0) -- (6.0,0);
    \draw[line width=1.5pt, hscink] (0,3.0) -- (6.0,3.0);
    \draw[line width=1.2pt, hscink] (0,0) -- (0,1.0);
    \draw[line width=1.2pt, hscink] (0,2.0) -- (0,3.0);
    \draw[line width=1.2pt, hscink, fill=white] (-0.16,1.0) rectangle (0.16,2.0);
    \node[font=\small, anchor=west] at (0.2,1.5) {$R$};
    \draw[line width=1pt, hscmuted, dashed] (3.0,-0.2) -- (3.0,3.2);
    \draw[line width=3.2pt, hscmuted] (3.8,-0.2) -- (3.8,3.2);
    \draw[velocity] (3.92,2.0) -- (4.9,2.0);
    \node[text=hscvel, font=\small, fill=white, inner sep=1pt, anchor=west]
      at (4.92,2.0) {$v$};
    \draw[force] (3.68,1.0) -- (2.7,1.0);
    \node[font=\footnotesize, text=hscforce, fill=white, inner sep=1pt, anchor=east]
      at (2.65,1.0) {$F = lIB$};
    \draw[->, line width=1pt, hscink] (4.05,0.2) -- (4.05,0.95)
      node[right, font=\small] {$I$};
    \draw[->, line width=1pt, hscink] (2.2,3.28) -- (1.2,3.28) node[left, font=\small] {$I$};
    \draw[->, line width=1pt, hscink] (1.2,-0.28) -- (2.2,-0.28) node[right, font=\small] {$I$};
    \draw[->, line width=1pt, hscink] (-0.45,2.0) -- (-0.45,1.0)
      node[midway, left, font=\small] {$I$};
    \draw[<->, line width=0.6pt, hscmuted] (6.35,0) -- (6.35,3.0)
      node[midway, right, text=hscink, font=\small] {$l$};
    \draw[<->, line width=0.6pt, hscmuted] (3.0,-0.62) -- (3.8,-0.62);
    \node[font=\small, below] at (3.4,-0.64) {$v\Delta t$};
  \end{tikzpicture}
  \captionof{figure}{The rod on rails. It sweeps the shaded area in time
    $\Delta t$.}\label{ch3:fig:rails}
\nextside
  \begin{itemize}
    \item On rails (\cref{ch3:fig:rails}) the current flows up the rod: anticlockwise.
    \item The force $F = lIB$ on the rod \emph{opposes} the motion (Lenz). The push
          does work, which becomes heat in $R$.
    \item Faraday's view: the rod sweeps out area, so the flux grows. Both views agree.
  \end{itemize}
\end{sidebyside}

\begin{itemize}
  \item Rod moving along the field lines, or along its own length: $\varepsilon = 0$.
  \item Induction: motion causes a current. Motor effect: a current causes motion.
\end{itemize}

\begin{formulacard}[notonsheet]{Motional emf}
  \cardformula{\varepsilon = lvB}
  \begin{symbols}
    \sym{$\varepsilon$}{emf between the ends of the rod}{\si{V}}
    \sym{$l$}{length of the rod inside the field}{\si{m}}
    \sym{$v$}{speed of the rod across the field}{\si{m.s^{-1}}}
    \sym{$B$}{magnetic field strength}{\si{T}}
  \end{symbols}
  \whentouse{rod, velocity and field all at \ang{90} to each other.}
  \derive{in time $\Delta t$ the rod sweeps an area $\Delta A = l\,v\Delta t$
    (\cref{ch3:fig:rails}), so $\Delta\Phi = B\,l\,v\Delta t$. Faraday's law with
    $N = 1$: $\lvert\varepsilon\rvert = \dfrac{\Delta\Phi}{\Delta t} = lvB$.}
  \trap{$l$ is the length \emph{inside the field}, not the whole wire.}
\end{formulacard}

%% ===================================================================
\section{Solenoids, metal plates and two coils}\label{ch3:sec:solenoids}
\covers{M6.IQ3.02}

Any change in flux induces an emf, whatever causes it. Magnets in a
\term{solenoid} (a long coil), metal plates and nearby coils follow the same laws.

\subsection{A magnet moving through a solenoid}
\begin{itemize}
  \item Approaching: the flux increases, so an emf is induced one way. Centre: the flux
        is largest and not changing, so $\varepsilon = 0$. Leaving: the flux decreases,
        so the emf is the other way (\cref{ch3:fig:falling}).
\end{itemize}

\begin{sidebyside}[0.64]
  \centering
  \includegraphics[width=\linewidth]{ch3_falling.pdf}
  \captionof{figure}{Model: a magnet falling through a coil.}\label{ch3:fig:falling}
\nextside
  \begin{itemize}[leftmargin=1.1em]
    \item A faster magnet gives taller, narrower pulses. A falling magnet leaves
          faster, so pulse 2 is taller and narrower.
    \item Each pulse has the same area, $N\lvert\Delta\Phi\rvert$.
  \end{itemize}
\end{sidebyside}

\begin{itemize}
  \item Closed circuit: the coil repels the magnet above it and attracts it below.
        Some gravitational potential energy becomes heat, so the magnet falls slower.
\end{itemize}

\subsection{Metal plates: eddy currents}
\begin{itemize}
  \item \term{Eddy currents}: loops of current induced inside solid metal, only where
        the flux through it changes. Lenz's law gives their direction.
  \item Plate moving right, field into the page: anticlockwise where the plate enters
        the field; none when fully inside; clockwise where it leaves.
  \item The force on each eddy loop opposes the motion. The plate slows, and its
        kinetic energy becomes heat: \term{magnetic braking}. Uses: brakes, induction
        cooktops, metal detectors. Problem: energy loss in transformer cores.
\end{itemize}

\subsection{One solenoid near another}
\begin{itemize}
  \item \term{Mutual induction}: a changing current in the \term{primary coil} (joined
        to the supply) induces an emf in a nearby \term{secondary coil}. A coil with a
        steady current, moved in or out, acts like a bar magnet.
  \item Switch closed: a brief emf; the secondary's induced field opposes the growing
        field. Kept closed: no emf. Opened: a brief emf the other way.
  \item A \term{soft iron} core (easily magnetised) guides more flux into the
        secondary: more emf.
  \item An \term{AC} supply (current keeps changing size and direction) gives an emf
        all the time: a transformer. A steady \term{DC} supply gives an emf only when
        switched.
  \item The secondary's energy comes from the supply (or your work, if a coil moves).
  \item \textbf{A coil wound on a long air-core solenoid} (all the primary's flux passes
        through it): each turn has $\Phi = BA$, with
        $B = \frac{\mu_0 N_\mathrm{p} I_\mathrm{p}}{L}$. Faraday's law then gives
        $\lvert\varepsilon_\mathrm{s}\rvert = N_\mathrm{s}\,\frac{\mu_0 N_\mathrm{p}
        A}{L}\,\frac{\lvert\Delta I_\mathrm{p}\rvert}{\Delta t}$ (not printed in this
        form). A faster change of $I_\mathrm{p}$ gives a bigger emf. $N_\mathrm{p}$,
        $N_\mathrm{s}$ = turns of the primary, of the secondary; $A$, $L$ = cross-section
        area and length of the primary.
\end{itemize}

%% ===================================================================
\section{Key experiments}\label{ch3:sec:experiments}
\covers{M6.IQ3.02}

\begin{experiment}{Faraday's induction experiments (1831)}
\aim{To show that a changing flux induces an emf.}
\method{A magnet was pushed into and pulled out of a coil joined to a galvanometer (a
  meter for very small currents). A wire was moved across a magnet's field. On two coils sharing an iron ring, a
  switch was closed and opened.}
\results{The needle flicked only while something changed. Opposite changes gave
  opposite flicks. Faster changes gave bigger flicks.}
\conclusion{An emf is induced only while the flux changes. Only relative motion
  matters.}
\significance{With Ørsted's work (1820), it linked electricity and magnetism.}
\end{experiment}

\begin{experiment}{Eddy currents: falling magnets and swinging plates}
\aim{To show that eddy currents oppose the motion that causes them.}
\method{A magnet fell through PVC, copper and slotted copper tubes. Solid and
  slotted plates swung between magnet poles.}
\results{Fall times: PVC \qty{0.41}{s} (free fall), slotted copper \qty{0.56}{s},
  copper \qty{3.57}{s}. The solid plate stopped fast; the slotted one swung on.}
\conclusion{Eddy currents in the copper oppose the fall (Lenz); the lost potential
  energy becomes heat. Slots break the current loops, so braking is weaker.}
\significance{Evidence for Lenz's law and energy conservation; used in magnetic
  brakes.}
\end{experiment}

%% ===================================================================
\begin{summarybox}{Remember}
\begin{itemize}
  \item $\Phi = BA\cos\theta$, $\theta$ from the normal. It changes if $B$, $A$ or
        $\theta$ changes.
  \item Only a \emph{changing} flux induces an emf:
        $\varepsilon = -N\dfrac{\Delta\Phi}{\Delta t}$.
  \item Lenz: the induced current opposes the \emph{change} in flux (energy is conserved).
  \item Moving rod: $\varepsilon = lvB$ (not on the sheet).
  \item Magnet through a coil: two opposite pulses. Two coils: an emf only while the
        primary current changes.
\end{itemize}
\end{summarybox}
}{}% Electromagnetic Induction: Flux, Faraday and Lenz
\IfFileExists{./ch4.tex}{\setcounter{chapter}{3}% ch4.tex -- Module 6 ESSENTIALS, Chapter 4: Transformers and Power Transmission
% Knowledge points only (no questions). Condensed from src/mod6/ch4.tex
% (final full book): same notation and facts. The two figures are TikZ, copied from the
% full chapter (shrunk). Dot points: M6.IQ3.03, M6.IQ3.04, M6.IQ3.05.
\chapter{Transformers and Power Transmission}\label{ch4:chap}
\covers{M6.IQ3.03, M6.IQ3.04, M6.IQ3.05}

% chapter-local colours (prefix chFour, as in the full book)
\colorlet{chFourcopper}{hscseriesB!80!black}  % coil wire
\colorlet{chFouriron}{hscmuted!28}            % iron core
\colorlet{chFourcurrent}{hscseriesB}          % induced (eddy) currents

\begin{summarybox}{In this chapter}
\begin{itemize}
  \item How a transformer changes an AC voltage, and the ideal-transformer equations.
  \item Why real transformers waste energy, and how engineers reduce the losses.
  \item Why the grid uses high voltage, and other uses of transformers.
\end{itemize}
\end{summarybox}

%% ==================================================================
\section{The ideal transformer}\label{ch4:sec:ideal}
\covers{M6.IQ3.03}

\begin{bigidea}
A transformer changes the size of an AC voltage. A changing current in one coil makes a
changing flux in an iron core, and this induces an emf in a second coil. Energy is
conserved: voltage up means current down.
\end{bigidea}

\begin{keywords}
  \kw{Alternating current (AC)}{current that keeps reversing (DC flows one way only).}
  \kw{Primary / secondary coil}{the input coil (joined to the supply) / the output coil
    (joined to the load, the device that uses the energy).}
  \kw{Mutual induction}{an emf made in one coil by a changing current in a nearby coil.}
  \kw{Step-up / step-down}{more / fewer secondary turns: bigger / smaller voltage.}
\end{keywords}

\subsection{How a transformer works}

Both coils are wrapped around one closed \term{core}: a ring of \term{soft iron} (iron that
is easy to magnetise and demagnetise), as in \cref{ch4:fig:transformer}.

\textbf{How the output voltage is made (cause $\to$ effect):}
\begin{enumerate}
  \item AC in the primary makes a magnetic field that keeps changing size and direction.
  \item The iron core carries this changing flux to the secondary coil.
  \item The flux through each secondary turn changes, so an emf is induced in every turn
        (Faraday's law, $\varepsilon = -N\dfrac{\Delta\Phi}{\Delta t}$).
  \item The turns are in series, so their emfs add up: more turns, bigger voltage.
  \item The flux reverses with the supply, so the output is AC of the \emph{same frequency}.
\end{enumerate}

\begin{sidebyside}[0.58]
  \centering
  \begin{tikzpicture}[x=0.7cm, y=0.75cm]
    % iron core (closed loop)
    \fill[chFouriron, even odd rule] (0,0) rectangle (5.6,3.6) (1.1,0.9) rectangle (4.5,2.7);
    \draw[hscink!70, line width=0.6pt] (0,0) rectangle (5.6,3.6) (1.1,0.9) rectangle (4.5,2.7);
    % flux path (clockwise at this instant)
    \draw[hscfield, line width=0.8pt, dashed] (0.55,0.45) rectangle (5.05,3.15);
    \draw[field] (1.4,3.15) -- (4.2,3.15);
    \draw[field] (4.2,0.45) -- (1.4,0.45);
    \node[diagramlabel, text=hscfield, align=center] at (2.8,1.8)
      {changing\\ flux $\Phi$};
    % primary coil: 4 turns on the left limb
    \foreach \y in {1.2,1.6,2.0,2.4}
      \draw[chFourcopper, line width=1.6pt] (-0.14,\y) -- (1.24,\y+0.18);
    % secondary coil: 8 turns on the right limb
    \foreach \y in {0.95,1.17,...,2.5}
      \draw[chFourcopper, line width=1.6pt] (5.74,\y) -- (4.36,\y+0.11);
    % primary circuit and AC supply
    \draw[hscink, line width=0.7pt] (-0.14,2.4) -- (-2.2,2.4) -- (-2.2,2.2);
    \draw[hscink, line width=0.7pt] (-0.14,1.2) -- (-2.2,1.2) -- (-2.2,1.4);
    \draw[hscink, line width=0.7pt, fill=white] (-2.2,1.8) circle (0.25cm);
    \node at (-2.2,1.8) {$\sim$};
    \node[diagramlabel, align=center, below] at (-2.2,1.1) {AC\\ supply};
    \node[diagramlabel] at (-1.2,1.8) {$V_\mathrm{p}$};
    \node[diagramlabel, above] at (-1.2,2.4) {$I_\mathrm{p}$};
    % secondary circuit and lamp
    \draw[hscink, line width=0.7pt] (5.74,2.49) -- (7.6,2.49) -- (7.6,2.15);
    \draw[hscink, line width=0.7pt] (5.74,0.95) -- (7.6,0.95) -- (7.6,1.45);
    \draw[hscink, line width=0.7pt, fill=white] (7.6,1.8) circle (0.21cm);
    \draw[hscink, line width=0.7pt] (7.6,1.8) ++(45:0.21cm) -- ++(225:0.42cm)
      (7.6,1.8) ++(135:0.21cm) -- ++(315:0.42cm);
    \node[diagramlabel, align=center, below] at (7.6,0.85) {lamp\\ (load)};
    \node[diagramlabel] at (6.7,1.75) {$V_\mathrm{s}$};
    \node[diagramlabel, above] at (6.7,2.49) {$I_\mathrm{s}$};
    % labels
    \node[diagramlabel] at (-0.4,4.1) {primary coil ($N_\mathrm{p}$ turns)};
    \draw[hscmuted, line width=0.5pt] (0.3,3.8) -- (0.3,2.52);
    \node[diagramlabel] at (6.0,4.1) {secondary coil ($N_\mathrm{s}$ turns)};
    \draw[hscmuted, line width=0.5pt] (5.3,3.8) -- (5.3,2.58);
    \node[diagramlabel, below] at (2.8,-0.05) {soft-iron core (a closed loop)};
  \end{tikzpicture}
  \captionof{figure}{A step-up transformer ($N_\mathrm{s} > N_\mathrm{p}$). The flux arrows
    show one instant; they reverse every half-cycle. No wire joins the coils: no charge
    crosses, and the changing field carries the energy.}
  \label{ch4:fig:transformer}
\nextside
  \begin{itemize}
    \item \textbf{Why AC is needed.} Steady DC makes a \emph{steady} field, so
          $\Delta\Phi/\Delta t = 0$: no emf. DC gives only a short pulse when it is
          switched on or off.
    \item \textbf{Why $V \propto N$.} The same flux passes through every turn of both
          coils, so each turn has the same emf. So the voltage ratio equals the
          \term{turns ratio} $N_\mathrm{p} : N_\mathrm{s}$.
  \end{itemize}
\end{sidebyside}

\begin{formulacard}{The transformer (turns) equation}
  \cardformula{\frac{V_\mathrm{p}}{V_\mathrm{s}} = \frac{N_\mathrm{p}}{N_\mathrm{s}}}
  \begin{symbols}
    \sym{$V_\mathrm{p}$, $V_\mathrm{s}$}{voltage of the primary (input), of the secondary}{\si{V}}
    \sym{$N_\mathrm{p}$, $N_\mathrm{s}$}{number of turns on the primary, on the secondary}{no unit}
  \end{symbols}
  \whentouse{a transformer on AC (exact if ideal). Not for DC, not for currents.}
  \trap{primary values on top of both fractions. The primary is joined to the supply.}
\end{formulacard}

\textbf{What the output looks like.}
\begin{itemize}
  \item The frequency stays the same. The \term{peak} (largest) voltage is multiplied by
        $N_\mathrm{s}/N_\mathrm{p}$.
  \item Mains: \qty{230}{V} AC, \qty{50}{Hz} ($T = 1/f = \qty{20}{ms}$). \qty{230}{V} is an
        \term{effective value} (the steady voltage that heats a heater by the same amount);
        the peak is about \qty{325}{V}. Never mix peak and effective values.
  \item \textbf{A graph of the primary \emph{current}:} the flux follows the current, so
        the secondary emf follows the \term{gradient} (steepness) of the current graph.
        Flat: zero emf. Steepest: biggest emf. A straight-line rise: a constant emf. A sine
        current: a cosine-shaped emf, shifted by a quarter of a period.
\end{itemize}

\subsection{Power and current}

In an \term{ideal transformer} no energy is wasted, so power in equals power out. The
\term{power} (energy each second) of each coil is $P = VI$ (on the sheet).

\begin{formulacard}{Power in an ideal transformer}
  \cardformula{V_\mathrm{p} I_\mathrm{p} = V_\mathrm{s} I_\mathrm{s}}
  \begin{symbols}
    \sym{$V_\mathrm{p}$, $I_\mathrm{p}$}{voltage and current of the primary (input)}{\si{V}, \si{A}}
    \sym{$V_\mathrm{s}$, $I_\mathrm{s}$}{voltage and current of the secondary (output)}{\si{V}, \si{A}}
  \end{symbols}
  \whentouse{an ideal transformer, to find a current or a power.}
  \trap{step-up: \emph{voltage} up, current \emph{down}. Never multiply $V_\mathrm{p}$ by
    $I_\mathrm{s}$.}
\end{formulacard}

\begin{itemize}
  \item With the turns equation:
        $\dfrac{I_\mathrm{s}}{I_\mathrm{p}} = \dfrac{V_\mathrm{p}}{V_\mathrm{s}}
        = \dfrac{N_\mathrm{p}}{N_\mathrm{s}}$ (not printed in this form). The current ratio
        is upside down: step-down gives a \emph{bigger} secondary current.
  \item \textbf{The load decides the current.} The load takes the current it needs; the
        primary then draws just enough to supply that power.
\end{itemize}

%% ==================================================================
\section{Real transformers: energy losses and efficiency}\label{ch4:sec:losses}
\covers{M6.IQ3.04}

\begin{bigidea}
Real transformers waste a little energy, mostly as heat. So the power out is a bit less
than the power in. A closed iron core, thick wire and a laminated core keep the losses
small.
\end{bigidea}

\begin{keywords}
  \kw{Incomplete flux linkage}{some flux from the primary misses the secondary.}
  \kw{Eddy currents}{loops of current induced inside a solid piece of metal by a changing
    magnetic flux.}
  \kw{Laminated core}{a core made of thin iron sheets with insulation between them.}
\end{keywords}

Real: $V_\mathrm{s}I_\mathrm{s} < V_\mathrm{p}I_\mathrm{p}$. The lost energy becomes heat
(and a little sound). \Cref{ch4:tab:losses} gives the three syllabus limitations.

\begin{table}
  \caption{The three main limitations of the ideal transformer model.}\label{ch4:tab:losses}
  \centering\small
  \begin{tabular}{@{}L{25mm}L{64mm}L{48mm}@{}}
    \toprule
    \colhead{Limitation} & \colhead{Cause and effect} & \colhead{How to reduce it}\\
    \midrule
    Incomplete flux linkage & Some primary flux leaks into the air and misses the
      secondary, so the secondary emf is smaller than predicted. &
      One \emph{closed} loop of soft iron for both coils. Wind the secondary directly over
      the primary.\\
    \addlinespace
    Resistive heating in the coils & The current heats the long copper coils ($I^2R$).
      This \term{resistive heating} is heat made when a current flows through a
      resistance. &
      Thick copper wire (thickest on the coil with the larger current). Cooling (vents,
      fans or oil) keeps the wire cool, and a cool wire has less resistance.\\
    \addlinespace
    Eddy currents in the core & The changing flux induces loops of current in the iron
      core, which heat it. & A laminated core, or a ferrite core.\\
    \bottomrule
  \end{tabular}
\end{table}

\textbf{Evaluating the ideal model.} Big transformers in power stations and
substations are about 99\% efficient; small ones often waste 10--20\%. So the ideal model
is a very good approximation, but never exact. Losses can be reduced, never removed.

\begin{formulacard}[notonsheet]{Efficiency}
  \cardformula{\eta = \frac{P_\mathrm{out}}{P_\mathrm{in}} \times 100\%
    = \frac{V_\mathrm{s}I_\mathrm{s}}{V_\mathrm{p}I_\mathrm{p}} \times 100\%}
  \begin{symbols}
    \sym{$\eta$}{efficiency (Greek ``eta'')}{no unit (\%)}
    \sym{$P_\mathrm{out}$, $P_\mathrm{in}$}{useful power out, power in}{\si{W}}
  \end{symbols}
  \whentouse{a real transformer (ideal: $\eta = 100\%$). Energy wasted:
    $\Delta E = (P_\mathrm{in} - P_\mathrm{out})\,\Delta t$ (from $P = \Delta E/\Delta t$).}
  \derive{a definition: useful power out $\div$ power in.}
  \trap{the efficiency goes on the \emph{power}, not on the voltage. Use $0.92$, not $92$,
    inside a formula.}
\end{formulacard}

\subsection{Eddy currents and the laminated core}

\begin{itemize}
  \item The core is a conductor with a changing flux. So eddy currents flow in it, in loops
        at right angles to the flux. By Lenz's law they oppose the change in flux. They heat the core.
  \item A laminated core is thin iron sheets, each with a thin insulating layer. The sheets
        lie \emph{along} the flux, so the flux passes easily. The insulation cuts the
        eddy-current paths: each loop stays inside one sheet (\cref{ch4:fig:eddy}).
\end{itemize}

\begin{sidebyside}[0.6]
  \centering
  \begin{tikzpicture}[x=0.75cm, y=0.75cm]
    \tikzset{chFourdot/.pic={\draw[hscfield, line width=0.6pt, fill=white] (0,0) circle (0.12);
      \fill[hscfield] (0,0) circle (0.035);}}
    % (a) solid core
    \fill[chFouriron] (0,0) rectangle (3,3);
    \draw[hscink!70, line width=0.6pt] (0,0) rectangle (3,3);
    \draw[chFourcurrent, line width=1.3pt, rounded corners=7pt] (0.5,0.5) rectangle (2.5,2.5);
    \draw[chFourcurrent, line width=1.3pt, -{Stealth[length=2.6mm]}] (1.3,2.5) -- (1.75,2.5);
    \draw[chFourcurrent, line width=1.3pt, -{Stealth[length=2.6mm]}] (1.7,0.5) -- (1.25,0.5);
    \draw[chFourcurrent, line width=1.3pt, -{Stealth[length=2.6mm]}] (2.5,1.7) -- (2.5,1.25);
    \draw[chFourcurrent, line width=1.3pt, -{Stealth[length=2.6mm]}] (0.5,1.3) -- (0.5,1.75);
    \foreach \p in {(1.5,1.5),(0.25,0.25),(2.75,0.25),(0.25,2.75),(2.75,2.75)}
      \pic at \p {chFourdot};
    \node[font=\small\bfseries, text=hscink] at (1.5,3.5) {(a) solid core};
    % (b) laminated core
    \begin{scope}[xshift=2.9cm]
      \fill[chFouriron] (0,0) rectangle (3,3);
      \foreach \x in {0.5,1.0,...,2.6}
        \draw[hscink, line width=1.3pt] (\x,0) -- (\x,3);
      \draw[hscink!70, line width=0.6pt] (0,0) rectangle (3,3);
      \foreach \x in {0,0.5,...,2.6} {
        \draw[chFourcurrent, line width=1pt, rounded corners=3pt]
          (\x+0.13,0.35) rectangle (\x+0.37,2.65);
        \draw[chFourcurrent, line width=1pt, -{Stealth[length=2mm]}]
          (\x+0.37,1.7) -- (\x+0.37,1.3);
        \draw[chFourcurrent, line width=1pt, -{Stealth[length=2mm]}]
          (\x+0.13,1.3) -- (\x+0.13,1.7);
      }
      \node[font=\small\bfseries, text=hscink] at (1.5,3.5) {(b) laminated core};
      \node[diagramlabel, anchor=west] (ins) at (3.3,2.4) {insulating layer};
      \draw[hscmuted, line width=0.4pt] (ins.west) -- (2.5,2.2);
      \node[diagramlabel, anchor=west] (sh) at (3.3,1.5) {thin iron sheet};
      \draw[hscmuted, line width=0.4pt] (sh.west) -- (2.8,1.5);
      \node[diagramlabel, anchor=west] (ed) at (3.3,0.6) {eddy current};
      \draw[hscmuted, line width=0.4pt] (ed.west) -- (2.87,0.6);
    \end{scope}
    \pic at (-0.45,-0.6) {chFourdot};
    \node[diagramlabel, anchor=west] at (-0.25,-0.6)
      {flux $\Phi$ points out of the page in both cores};
  \end{tikzpicture}
  \captionof{figure}{A slice through the core: (a) large eddy currents, (b) small ones. The
    flux comes out of the page and is growing. So the eddy currents flow clockwise (they
    reverse every half-cycle).}
  \label{ch4:fig:eddy}
\nextside
  \begin{itemize}
    \item \textbf{Why this works:} each thin loop surrounds only a thin strip of flux
          (small emf) and has a narrow, high-resistance path. So the eddy currents and
          their heat are much smaller.
    \item Small chargers use \term{ferrite} cores. Ferrite is a magnetic ceramic that
          conducts very poorly. Cooling does
          \emph{not} stop eddy currents.
  \end{itemize}
\end{sidebyside}


%% ==================================================================
\section{Sending electricity at high voltage}\label{ch4:sec:transmission}
\covers{M6.IQ3.05}

\begin{bigidea}
Power lines have resistance, so their current makes heat: $I^2R$. Stepping the voltage
up makes the current small, so much less energy is wasted.
\end{bigidea}


\begin{keywords}
  \kw{Transmission line}{a long, thick wire on tall towers that carries electricity far.}
  \kw{Voltage drop}{the voltage used up pushing current through a resistance, $IR$.}
  \kw{Substation}{a place with transformers that change the voltage.}
\end{keywords}

A long line has several ohms of resistance $R$. The current $I$ is the same all along it,
so the town gets only $V - IR$ (Ohm's law, \cref{ch4:fig:line}).

\begin{figure}
  \centering
  \begin{tikzpicture}[x=1cm, y=1cm]
    \draw[hscink, line width=0.7pt] (0,0.95) -- (0,1.6) -- (2.6,1.6);
    \draw[hscink, line width=0.7pt] (4.4,1.6) -- (7.4,1.6) -- (7.4,1.1);
    \draw[hscink, line width=0.7pt] (7.4,0.1) -- (7.4,-0.4) -- (0,-0.4) -- (0,0.25);
    \draw[hscink, line width=0.7pt, fill=white] (0,0.6) circle (0.35);
    \node at (0,0.6) {$\sim$};
    \draw[hscink, line width=0.7pt, fill=hscforce!10] (2.6,1.4) rectangle (4.4,1.8);
    \node[diagramlabel] at (3.5,1.6) {$R$};
    \draw[hscink, line width=0.7pt, fill=hscbrand!8] (7.0,0.1) rectangle (7.8,1.1);
    \node[diagramlabel, rotate=90] at (7.4,0.6) {town};
    \draw[velocity] (1.0,1.6) -- (1.9,1.6) node[midway, above, text=hscvel] {$I$};
    \node[diagramlabel, align=center, anchor=east] at (-0.45,0.6)
      {power station\\ sends voltage $V$};
    \node[diagramlabel, align=center, above] at (3.5,1.85)
      {line resistance $R$ (both wires)\\ voltage drop $V_\mathrm{drop} = IR$};
    \node[diagramlabel, align=center, anchor=west] at (7.95,0.6)
      {voltage left\\ for the town\\ $= V - IR$};
    \node[diagramlabel, text=hscforce] at (3.5,1.0) {heat: $P_\mathrm{loss} = I^2R$};
  \end{tikzpicture}
  \caption{A simple model of a transmission line. The current $I$ is the same all the way
    around the loop.}
  \label{ch4:fig:line}
\end{figure}

\textbf{Why high voltage (cause $\to$ effect):}
\begin{enumerate}
  \item A power station sends a fixed power. Since $P = VI$, a higher voltage needs a
        smaller current.
  \item The loss $I^2R$ depends on the \emph{square} of the current: 10 times the voltage
        gives $\tfrac{1}{10}$ of the current and $\tfrac{1}{100}$ of the loss. $R$ stays the
        same. (Thicker wires also help, but they are heavy and very expensive.)
\end{enumerate}

\begin{formulacard}[notonsheet, unbreakable]{Power lost in a transmission line}
  \cardformula{P_\mathrm{loss} = I^2R}
  \begin{symbols}
    \sym{$P_\mathrm{loss}$}{power turned into heat in the line}{\si{W}}
    \sym{$I$}{current in the line}{\si{A}}
    \sym{$R$}{total resistance of the line (both wires)}{\si{\ohm}}
  \end{symbols}
  \whentouse{heat lost in lines or coils. Find $I = P/V$ first; the town gets
    $P - P_\mathrm{loss}$.}
  \derive{$P = VI$ for the line only, with $V_\mathrm{drop} = IR$ (both on the sheet):
    $P_\mathrm{loss} = V_\mathrm{drop}\,I = (IR)\,I = I^2R$.}
  \trap{do not use $\dfrac{V^2}{R}$ with the line's voltage (for example \qty{330}{kV}).
    There, $V$ must be the voltage drop $IR$, not the voltage sent.}
\end{formulacard}

\textbf{The grid.} The \term{electricity grid} is the network of power stations, lines and transformers.
Typical NSW voltages: power station about \qty{20}{kV} $\to$ step-up $\to$ transmission
lines \qty{330}{kV} or \qty{500}{kV} $\to$ step-down at substations $\to$ \qty{132}{kV} or
\qty{66}{kV} $\to$ \qty{11}{kV} $\to$ step-down on a power pole $\to$ homes \qty{230}{V}.
\begin{itemize}
  \item \textbf{Step up} at the power station: a low current in the long lines, so little
        energy is lost.
  \item \textbf{Step down} near towns: high voltages are dangerous. They can spark (arc)
        across air and need tall towers and big insulators. Appliances are built for
        \qty{230}{V}.
  \item \textbf{Why AC.} Transformers need a changing current: AC, not steady DC. (A few
        undersea links, such as Basslink, use DC.)
  \item \textbf{Other uses.} Step-up: microwave oven (\qty{230}{V} $\to$ about \qty{2}{kV}).
        Step-down: welder (\qty{230}{V} $\to$ \qty{50}{V}: a very large current melts
        metal); phone charger (\qty{230}{V} $\to$ \qty{5}{V}, small ferrite core, then
        electronics make DC).
\end{itemize}


%% ==================================================================
\section{Key experiments}\label{ch4:sec:experiments}
\covers{M6.IQ3.03, M6.IQ3.04}

\begin{experiment}{Testing the transformer equations}
\aim{To test $V_\mathrm{p}/V_\mathrm{s} = N_\mathrm{p}/N_\mathrm{s}$ for a real
  transformer, and to find its efficiency.}
\method{A demountable (take-apart) kit: laminated U-shaped iron core, iron \emph{I-bar} to close the
  top, slide-on coils; low-voltage AC, AC meters. Primary 600 turns at \qty{6.0}{V}:
  measure $V_\mathrm{s}$ for 150 to 1200 secondary turns. With a load, read all four
  meters. Then lift off the I-bar.}
\results{$V_\mathrm{s}$ against $N_\mathrm{s}$: a straight line through the origin, about
  5\% below the ideal line. With the load, $\eta \approx 90\%$. Opening the core
  cut $V_\mathrm{s}$ (300 turns) from \qty{2.86}{V} to about \qty{0.9}{V}.}
\conclusion{$V_\mathrm{s} \propto N_\mathrm{s}$, as predicted. The small gap is mainly flux
  leakage; $P_\mathrm{out} < P_\mathrm{in}$ (heat). The ideal model is close, not exact.}
\significance{It supports the model behind every grid transformer, and shows why cores are
  closed loops.}
\end{experiment}

%% ==================================================================
\begin{summarybox}{Remember}
\begin{itemize}
  \item A transformer needs AC: changing primary current $\to$ changing flux in the iron
        core $\to$ induced emf in the secondary (Faraday's law). Steady DC gives no output.
  \item $\dfrac{V_\mathrm{p}}{V_\mathrm{s}} = \dfrac{N_\mathrm{p}}{N_\mathrm{s}}$: more
        secondary turns, bigger voltage. The frequency stays the same.
  \item Ideal: $V_\mathrm{p}I_\mathrm{p} = V_\mathrm{s}I_\mathrm{s}$, so voltage up means
        current down: $\dfrac{I_\mathrm{s}}{I_\mathrm{p}} = \dfrac{N_\mathrm{p}}{N_\mathrm{s}}$.
  \item Real losses: flux leakage, $I^2R$ heating in the coils, eddy currents in the core.
        Fixes: closed soft-iron core, thick copper wire, laminated core.
  \item Not on the sheet: $\eta = \dfrac{P_\mathrm{out}}{P_\mathrm{in}} \times 100\%$ and
        $P_\mathrm{loss} = I^2R$.
  \item Fixed power: higher voltage $\to$ smaller current $\to$ much less $I^2R$ loss.
        Step up at the power station; step down in stages to \qty{230}{V} for safety.
\end{itemize}
\end{summarybox}
}{}% Transformers and Power Transmission
\IfFileExists{./ch5.tex}{\setcounter{chapter}{4}% ch5.tex -- Module 6 ESSENTIALS, Chapter 5: Applications of the motor effect
% Condensed from src/mod6/ch5.tex: knowledge points only.
\chapter[Motors, Generators and Braking]{Applications of the Motor Effect: Motors, Generators and Braking}\label{ch5:chap}
\covers{M6.IQ4.01, M6.IQ4.02, M6.IQ4.03}

% ---- chapter-local TikZ helpers (prefix chfive, as in the full book) ----
\tikzset{
  chfiveIn/.pic={\draw[hscink, line width=0.7pt, fill=white] (0,0) circle (0.16);
    \draw[hscink, line width=0.7pt] (-0.11,-0.11) -- (0.11,0.11) (-0.11,0.11) -- (0.11,-0.11);},
  chfiveOut/.pic={\draw[hscink, line width=0.7pt, fill=white] (0,0) circle (0.16);
    \fill[hscink] (0,0) circle (0.05);},
  chfiveWire/.style={hscink, line width=1.4pt},
  pics/chfivePoles/.style={code={
    \fill[hscforce!15, draw=hscink, line width=0.5pt] (-2.05,-1.15) rectangle (-1.8,1.15);
    \node[font=\small\bfseries] at (-1.925,1.38) {N};
    \fill[hscvel!15, draw=hscink, line width=0.5pt] (1.8,-1.15) rectangle (2.05,1.15);
    \node[font=\small\bfseries] at (1.925,1.38) {S};
    \foreach \yy in {-0.95,-0.35,0.35,0.95}
      \draw[field, hscfield!45] (-1.8,\yy) -- (1.8,\yy);}},
}

\begin{summarybox}{In this chapter}
\begin{itemize}
  \item The torque on a coil, and the job of each part of a DC motor.
  \item Back emf: the large start-up current and the top speed of a motor.
  \item AC and DC generators, AC induction motors and magnetic brakes.
  \item Lenz's law as conservation of energy in all of these devices.
\end{itemize}
\end{summarybox}

%% ------------------------------------------------------------------
\section{Torque on a coil and the simple DC motor}\label{ch5:sec:torque}
\covers{M6.IQ4.01}

\begin{bigidea}
A current-carrying coil in a magnetic field feels opposite forces on its two sides, so
it turns. The torque is largest when its plane is parallel to the field.
\end{bigidea}

\begin{keywords}
  \kw{Normal}{a line at \ang{90} to the flat face (the plane) of the coil.}
  \kw{Torque ($\tau$)}{the turning effect of a force. Unit: newton metre (\si{N.m}).}
  \kw{Stator}{the part of a motor that does not move. It makes the magnetic field.}
  \kw{Rotor (armature)}{the part that turns: a coil on a soft-iron core, fixed to the axle.}
  \kw{Split-ring commutator}{two half-rings that reverse the current every half turn.}
  \kw{Brushes}{fixed carbon blocks that connect the spinning coil to the supply.}
\end{keywords}

\begin{itemize}
  \item A \term{coil} is many loops (turns) of wire. Each long side is at \ang{90} to
        $B$. With $n$ turns, it feels a force $F = nlIB$.
  \item The long sides carry opposite currents, so one force is up and one is down. They
        form a \term{couple}: equal, opposite forces, not in line, that turn an object.
  \item The forces on the short end sides point along the axle. They give no torque.
  \item The forces keep the same size. Only the \term{lever arm} $r_\perp$ (axle to the
        line of the force) changes as the coil turns (\cref{ch5:fig:endon}).
  \item Only one side in the field: the force on that side is $F = lIB$ (the side is at
        \ang{90} to $B$), and $\tau = r_\perp F$.
\end{itemize}

\begin{formulacard}{Torque on a coil}
  \cardformula{\tau = nIA_\perp B = nIAB\sin\theta}
  \begin{symbols}
    \sym{$\tau$}{torque on the coil}{\si{N.m}}
    \sym{$n$}{number of turns in the coil}{no unit}
    \sym{$I$}{current in the coil}{\si{A}}
    \sym{$A$, $A_\perp$}{area of one turn of the coil; $A_\perp = A\sin\theta$}{\si{m^2}}
    \sym{$B$}{magnetic field strength}{\si{T}}
    \sym{$\theta$}{angle between $B$ and the normal to the coil}{degree (\si{\degree})}
  \end{symbols}
  \whentouse{a flat coil turning on an axle in a uniform field.}
  \trap{$\theta$ is from the \emph{normal}: plane at \ang{30} to $B$ means $\theta = \ang{60}$.}
\end{formulacard}

\begin{figure}
  \centering
  \begin{tikzpicture}[x=1cm, y=1cm]
    % (a) theta = 90
    \begin{scope}
      \pic {chfivePoles};
      \draw[guide, ->] (1.06,0) -- (1.62,0) node[pos=0.45, below, inner sep=1pt, font=\footnotesize, text=hscmuted] {$B$};
      \draw[->, hscmuted, line width=0.7pt, dashed] (0,0) -- (90:1.25) node[right, font=\footnotesize, text=hscmuted] {normal};
      \draw[hscmuted] (0.42,0) arc (0:90:0.42);
      \node[font=\footnotesize, text=hscmuted] at (45:0.62) {$\theta$};
      \draw[hscink, line width=2.2pt] (180:0.85) -- (0:0.85);
      \pic at (180:0.85) {chfiveOut};
      \pic at (0:0.85) {chfiveIn};
      \draw[force] ($(180:0.85)+(0,0.18)$) -- ++(0,0.62) node[left, text=hscforce, font=\footnotesize] {$F$};
      \draw[force] ($(0:0.85)+(0,-0.18)$) -- ++(0,-0.62) node[right, text=hscforce, font=\footnotesize] {$F$};
      \fill[hscink] (0,0) circle (1.6pt);
      \node[font=\footnotesize, align=center, anchor=north] at (0,-1.7)
        {(a) Plane $\parallel B$, $\theta = \ang{90}$:\\largest torque.};
    \end{scope}
    % (b) theta = 50
    \begin{scope}[xshift=4.75cm]
      \pic {chfivePoles};
      \draw[guide, ->] (0,0) -- (1.3,0) node[pos=0.85, below, inner sep=1pt, font=\footnotesize, text=hscmuted] {$B$};
      \draw[->, hscmuted, line width=0.7pt, dashed] (0,0) -- (50:1.25) node[above, font=\footnotesize, text=hscmuted] {normal};
      \draw[hscmuted] (0.42,0) arc (0:50:0.42);
      \node[font=\footnotesize, text=hscmuted] at (25:0.62) {$\theta$};
      \draw[hscink, line width=2.2pt] (140:0.85) -- (-40:0.85);
      \draw[guide] (140:0.85) -- ({0.85*cos(140)},0);
      \draw[hscacc, line width=1.4pt] ({0.85*cos(140)},0) -- (0,0);
      \node[font=\footnotesize, text=hscacc, below] at ({0.425*cos(140)},0) {$r_\perp$};
      \pic at (140:0.85) {chfiveOut};
      \pic at (-40:0.85) {chfiveIn};
      \draw[force] ($(140:0.85)+(0,0.18)$) -- ++(0,0.62) node[left, text=hscforce, font=\footnotesize] {$F$};
      \draw[force] ($(-40:0.85)+(0,-0.18)$) -- ++(0,-0.62) node[right, text=hscforce, font=\footnotesize] {$F$};
      \fill[hscink] (0,0) circle (1.6pt);
      \node[font=\footnotesize, align=center, anchor=north] at (0,-1.7)
        {(b) $\theta = \ang{50}$: smaller lever\\arm, smaller torque.};
    \end{scope}
    % (c) theta = 0
    \begin{scope}[xshift=9.5cm]
      \pic {chfivePoles};
      \draw[->, hscmuted, line width=0.7pt, dashed] (0,0) -- (0:1.25) node[pos=0.62, above=0.5pt, inner sep=0.8pt, fill=white, font=\scriptsize, text=hscmuted] {normal};
      \draw[hscink, line width=2.2pt] (90:0.85) -- (-90:0.85);
      \pic at (90:0.85) {chfiveOut};
      \pic at (-90:0.85) {chfiveIn};
      \draw[force] ($(90:0.85)+(0,0.18)$) -- ++(0,0.42) node[right, text=hscforce, font=\footnotesize] {$F$};
      \draw[force] ($(-90:0.85)+(0,-0.18)$) -- ++(0,-0.42) node[right, text=hscforce, font=\footnotesize] {$F$};
      \fill[hscink] (0,0) circle (1.6pt);
      \node[font=\footnotesize, align=center, anchor=north] at (0,-1.7)
        {(c) Plane $\perp B$ ($\theta = 0$): forces\\in line, zero torque.};
    \end{scope}
  \end{tikzpicture}
  \caption{A coil seen end-on (axle = black dot). $\odot$ = current out of the page,
    $\otimes$ = into the page. $\theta$ is between $B$ and the normal.}\label{ch5:fig:endon}
\end{figure}


\subsection{The simple DC motor}
\begin{sidebyside}[0.49]
  \centering
  \begin{tikzpicture}[x={(1cm,0cm)}, y={(0cm,1cm)}, z={(0.42cm,0.30cm)}, scale=0.72]
    \def\mz{2.8}
    % N magnet (left)
    \filldraw[draw=hscink, line width=0.5pt, fill=hscforce!28] (-2.6,-1,0) -- (-2.6,-1,\mz) -- (-2.6,1,\mz) -- (-2.6,1,0) -- cycle;
    \filldraw[draw=hscink, line width=0.5pt, fill=hscforce!14] (-3.7,1,0) -- (-2.6,1,0) -- (-2.6,1,\mz) -- (-3.7,1,\mz) -- cycle;
    \filldraw[draw=hscink, line width=0.5pt, fill=hscforce!22] (-3.7,-1,0) -- (-2.6,-1,0) -- (-2.6,1,0) -- (-3.7,1,0) -- cycle;
    \node[font=\bfseries] at (-3.15,0,0) {N};
    % field lines
    \draw[field] (-2.6,0.62,1.4) -- (2.6,0.62,1.4);
    \draw[field] (-2.6,-0.92,0.25) -- (2.6,-0.92,0.25);
    % S magnet (right)
    \filldraw[draw=hscink, line width=0.5pt, fill=hscvel!14] (2.6,1,0) -- (3.7,1,0) -- (3.7,1,\mz) -- (2.6,1,\mz) -- cycle;
    \filldraw[draw=hscink, line width=0.5pt, fill=hscvel!28] (3.7,-1,0) -- (3.7,-1,\mz) -- (3.7,1,\mz) -- (3.7,1,0) -- cycle;
    \filldraw[draw=hscink, line width=0.5pt, fill=hscvel!22] (2.6,-1,0) -- (3.7,-1,0) -- (3.7,1,0) -- (2.6,1,0) -- cycle;
    \node[font=\bfseries] at (3.15,0,0) {S};
    % axle
    \draw[hscmuted, line width=0.9pt, dashed] (0,0,-1.9) -- (0,0,3.4);
    % coil: current enters the right side (into the page), leaves by the left side
    \draw[chfiveWire, midarrow] (1.4,0,0.3) -- (1.4,0,2.3);
    \draw[chfiveWire, midarrow] (1.4,0,2.3) -- (-1.4,0,2.3);
    \draw[chfiveWire, midarrow] (-1.4,0,2.3) -- (-1.4,0,0.3);
    \draw[chfiveWire] (1.4,0,0.3) -- (0.12,0,0.3) -- (0.12,0,-0.9);
    \draw[chfiveWire] (-1.4,0,0.3) -- (-0.12,0,0.3) -- (-0.12,0,-0.9);
    \node[font=\footnotesize, text=hscink] at (1.75,-0.2,1.6) {$I$};
    \node[font=\footnotesize, text=hscink] at (-1.62,0.25,1.6) {$I$};
    % forces
    \draw[force] (1.4,-0.08,1.3) -- (1.4,-1.0,1.3) node[right, text=hscforce, font=\footnotesize] {$F$};
    \draw[force] (-1.4,0.08,1.3) -- (-1.4,1.0,1.3) node[left, text=hscforce, font=\footnotesize] {$F$};
    \node[font=\footnotesize, text=hscfield] at (2.0,0.85,1.4) {$B$};
    % commutator and brushes at z = -0.9
    \begin{scope}[shift={(0,0,-0.9)}]
      \fill[white] (0,0) circle (0.36);
      \draw[hscseriesB, line width=3pt] (100:0.27) arc (100:260:0.27);
      \draw[hscseriesA, line width=3pt] (-80:0.27) arc (-80:80:0.27);
      \filldraw[fill=hscmuted!55, draw=hscink, line width=0.4pt] (-0.78,-0.09) rectangle (-0.34,0.09);
      \filldraw[fill=hscmuted!55, draw=hscink, line width=0.4pt] (0.34,-0.09) rectangle (0.78,0.09);
      \draw[hscink, line width=0.8pt] (-0.78,0) -- (-1.5,0) -- (-1.5,-2.0) -- (-0.36,-2.0);
      \draw[hscink, line width=0.8pt] (0.78,0) -- (1.5,0) -- (1.5,-2.0) -- (0.0,-2.0);
      \draw[hscink, line width=1.6pt] (-0.36,-2.14) -- (-0.36,-1.86);
      \draw[hscink, line width=0.9pt] (0,-2.3) -- (0,-1.7);
      \node[font=\footnotesize] at (0.22,-1.65) {$+$};
      \node[font=\footnotesize] at (-0.58,-1.65) {$-$};
      \node[diagramlabel, anchor=north] at (-0.18,-2.35) {DC supply};
      \draw[hscmuted, line width=0.4pt] (-0.24,-0.22) -- (-1.75,-0.95) node[anchor=north east, diagramlabel, inner sep=1pt, align=right] {split-ring\\commutator};
      \draw[hscmuted, line width=0.4pt] (0.65,-0.09) -- (1.75,-0.95) node[anchor=north west, diagramlabel, inner sep=1pt] {brush};
    \end{scope}
    \draw[hscmuted, line width=0.4pt] (0,0,3.4) -- ++(0.6,0.5) node[anchor=south west, diagramlabel, inner sep=1pt] {axle};
    \draw[hscmuted, line width=0.4pt] (-0.4,0,2.3) -- ++(0,1.05) node[anchor=south, diagramlabel, inner sep=1pt] {coil (rotor)};
    \node[diagramlabel, anchor=south, align=center] at (-3.15,1,2.9) {magnet\\(stator)};
  \end{tikzpicture}
  \captionof{figure}{A simple DC motor. Seen from the commutator end, the coil turns
    clockwise.}\label{ch5:fig:motor}
\nextside
  \begin{itemize}
    \item A \term{DC motor} changes electrical energy into rotational kinetic energy
          (\cref{ch5:fig:motor}).
    \item The split ring reverses the current every half turn, as the plane passes
          perpendicular to $B$. So the torque always turns the coil one way.
    \item At the swap the brushes touch the gaps. \term{Rotational inertia} (the
          tendency to keep spinning) carries the coil on.
  \end{itemize}
\end{sidebyside}

\textbf{No commutator:} the coil rocks, then stops with its plane perpendicular to $B$.
\textbf{Dead spot:} a single coil stopped there has no torque, so it cannot start by
itself. Many coils at different angles give a smooth torque.

\textbf{Radial field.} Curved \term{pole pieces} (the shaped ends of the magnets) and a
soft-iron core make a
\term{radial magnetic field}: its lines point straight out from the centre, like spokes.
The field lies along the plane of the coil, so $\theta = \ang{90}$ in every position. The
torque stays at $nIAB$ for almost the whole turn. The maximum is \emph{not} bigger; the
average is.

%% ------------------------------------------------------------------
\section{Back emf, maximum speed and energy}\label{ch5:sec:backemf}
\covers{M6.IQ4.01, M6.IQ4.03}

\begin{bigidea}
A spinning motor coil is also a generator. Its induced emf, the \term{back emf}
($\varepsilon_\mathrm{back}$), opposes the supply. The \term{net emf},
$V - \varepsilon_\mathrm{back}$, pushes the current.
\end{bigidea}

\begin{itemize}
  \item A spinning coil has a changing flux, so an emf is induced (Faraday's law). By
        Lenz's law it opposes the change that caused it, so it acts against the supply.
  \item The back emf is proportional to the speed. Every motor has one, AC or DC.
\end{itemize}

\begin{sidebyside}[0.40]
  \centering
  \begin{tikzpicture}
    \begin{axis}[hsc, scale only axis, width=4.0cm, height=2.2cm,
        xmin=0, xmax=5, ymin=0, ymax=1.25, xtick=\empty, ytick={0,1},
        yticklabels={0,$\dfrac{V}{R}$}, xlabel={time from switch-on}, ylabel={current}]
      \addplot[hscseriesB, domain=0:5, samples=100] {0.1+0.9*exp(-x)};
      \node[anchor=south west, font=\footnotesize, text=hscseriesB] at (axis cs:2.0,0.2) {levels off above 0};
      \node[anchor=west, font=\footnotesize, align=left] at (axis cs:0.45,0.95)
        {largest at switch-on};
    \end{axis}
  \end{tikzpicture}
  \captionof{figure}{Current in a DC motor after switch-on. The torque has the same
    shape.}\label{ch5:fig:startup}
\nextside
  \begin{enumerate}
    \item Switch-on: the coil is still, so $\varepsilon_\mathrm{back} = 0$. $R$ is
          small, so $I = \dfrac{V}{R}$ and the torque are large.
    \item Speeding up: the back emf grows, so the current and torque fall.
    \item \term{Maximum speed}: the torque just balances friction and the load. The
          back emf is a little less than $V$, so a small current still flows
          (\cref{ch5:fig:startup}). In an ideal motor (no friction, no load), the back
          emf $\approx V$ and the current $\approx 0$.
  \end{enumerate}
\end{sidebyside}

\begin{itemize}
  \item A bigger \term{load} (whatever the motor drives) slows the motor. The back emf
        falls, so the current and torque rise.
  \item A motor \term{stalls} when it is switched on but cannot turn. With no back emf,
        $I = \dfrac{V}{R}$ is large. The heating $I^2R$ can melt the insulation (``burn
        out''). A \term{starting resistor} in series limits the start-up current.
  \item Energy: $VI = \varepsilon_\mathrm{back}I + I^2R$. Power in $=$ mechanical power
        $+$ heat in the coil. So back emf is how the motor does work. A stalled motor turns
        all of its input energy into heat.
  \item \textbf{Lenz's law is conservation of energy.} Imagine the induced emf
        \emph{helped} the supply. A faster coil would get more current and torque, so it
        would speed up without limit. Energy would be made from nothing. So the induced
        emf must oppose.
\end{itemize}

\begin{formulacard}[notonsheet]{Current in a running motor}
  \cardformula{I = \frac{V - \varepsilon_\mathrm{back}}{R}}
  \begin{symbols}
    \sym{$I$}{current in the motor coil}{\si{A}}
    \sym{$V$}{supply emf (supply voltage)}{\si{V}}
    \sym{$\varepsilon_\mathrm{back}$}{back emf induced in the spinning coil}{\si{V}}
    \sym{$R$}{resistance of the coil}{\si{\ohm}}
  \end{symbols}
  \whentouse{a running DC motor. At switch-on or when stalled,
    $\varepsilon_\mathrm{back} = 0$, so $I = \dfrac{V}{R}$.}
  \derive{$V = IR$ (on the sheet) with the net emf: $V - \varepsilon_\mathrm{back} = IR$.
    Times $I$: $VI = \varepsilon_\mathrm{back}I + I^2R$ (power in $=$ mechanical $+$ heat).}
  \trap{the back emf is zero at switch-on, not at top speed.}
\end{formulacard}

%% ------------------------------------------------------------------
\section{Generators: AC and DC}\label{ch5:sec:generators}
\covers{M6.IQ4.02, M6.IQ4.03}

\begin{bigidea}
A \term{generator} changes kinetic energy into electrical energy. It is a motor run
backwards: an outside force turns the coil, the flux changes, and an emf is induced.
\end{bigidea}

\begin{itemize}
  \item As the coil turns, $\theta$ changes, so the flux $\Phi = BA\cos\theta$ changes
        and an emf is induced (Faraday's law: $\varepsilon = -N\dfrac{\Delta\Phi}{\Delta t}$).
  \item \textbf{Plane perpendicular to $B$:} flux largest ($BA$) but not changing for an
        instant, so the emf is \emph{zero}.
  \item \textbf{Plane parallel to $B$:} flux zero but changing fastest (the sides cut
        straight across the field lines), so the emf is \emph{largest}.
  \item \textbf{AC generator:} \term{slip rings} (two full rings, each always joined to
        the same end of the coil). The coil's emf reverses every half turn, so the
        output is \term{alternating current (AC)}: it keeps reversing direction.
  \item \textbf{DC generator:} a split-ring commutator swaps the connections every half
        turn, when the emf is zero. The output is \term{direct current (DC)}: one way
        only, in pulses (\cref{ch5:fig:gengraphs}).
\end{itemize}

\begin{figure}
  \centering
  \begin{tikzpicture}
    \begin{axis}[hsc, name=fa, axis x line=middle, scale only axis, width=10.2cm, height=1.0cm,
        xmin=0, xmax=2, ymin=-1.3, ymax=1.3, xtick={0,0.25,...,2}, xticklabels={},
        ytick={-1,0,1}, yticklabels={$-BA$,0,$BA$}, ylabel={$\Phi$},
        ylabel style={rotate=-90}]
      \addplot[hscseriesA, domain=0:2, samples=161] {cos(360*x)};
    \end{axis}
    \foreach \f/\ang in {0/90, 0.125/0, 0.25/90, 0.375/0, 0.5/90} {
      \coordinate (ic) at ($(fa.north west)!\f!(fa.north east)+(0,0.95)$);
      \draw[hscink, line width=1.6pt] ($(ic)+({\ang}:0.24)$) -- ($(ic)+({\ang+180}:0.24)$);
      \fill[hscink] (ic) circle (1pt);
    }
    \node[anchor=west, font=\footnotesize, text=hscmuted] at ($(fa.north west)!0.56!(fa.north east)+(0,0.95)$)
      {$\leftarrow$ the coil seen end-on (field $B$ to the right)};
    \node[anchor=south, font=\scriptsize, text=hscmuted, align=center, text depth=0.3em] at ($(fa.north west)!0.0!(fa.north east)+(0,1.22)$) {plane\\$\perp B$};
    \node[anchor=south, font=\scriptsize, text=hscmuted, align=center, text depth=0.3em] at ($(fa.north west)!0.125!(fa.north east)+(0,1.22)$) {plane\\$\parallel B$};
    \begin{axis}[hsc, name=fb, axis x line=middle, at={(fa.south west)}, anchor=north west, yshift=-0.3cm,
        scale only axis, width=10.2cm, height=1.0cm,
        xmin=0, xmax=2, ymin=-1.3, ymax=1.3, xtick={0,0.25,...,2}, xticklabels={},
        ytick={-1,0,1}, yticklabels={$-\varepsilon_\mathrm{peak}$,0,$\varepsilon_\mathrm{peak}$},
        ylabel={emf}]
      \addplot[hscseriesB, domain=0:2, samples=161] {sin(360*x)};
      \node[anchor=west, font=\footnotesize, align=left] at (axis cs:2.05,0) {AC\\(slip rings)};
    \end{axis}
    \begin{axis}[hsc, name=fc, at={(fb.south west)}, anchor=north west, yshift=-0.3cm,
        scale only axis, width=10.2cm, height=1.0cm,
        xmin=0, xmax=2, ymin=-0.3, ymax=1.3, xtick={0,0.25,...,2},
        xticklabels={0,$T/4$,$T/2$,$3T/4$,$T$,,,,$2T$},
        ytick={0,1}, yticklabels={0,$\varepsilon_\mathrm{peak}$}, ylabel={emf}, xlabel={time}]
      \addplot[hscseriesC, domain=0:2, samples=161] {abs(sin(360*x))};
      \node[anchor=west, font=\footnotesize, align=left] at (axis cs:2.05,0.5) {DC\\(split ring)};
    \end{axis}
  \end{tikzpicture}
  \caption{Flux and output emf for a coil that starts with its plane perpendicular to $B$
    and turns at a steady rate. $T$ is the time for one turn.}\label{ch5:fig:gengraphs}
\end{figure}

\begin{itemize}
  \item \textbf{Faster turning:} bigger peak emf \emph{and} higher frequency. Twice the
        speed gives twice the peak and half the period. More turns, a stronger field or
        a bigger coil raise the peak only. The other direction flips the graph.
  \item \textbf{Energy:} with a load connected, the current in the coil feels a
        motor-effect force. By Lenz's law it opposes the rotation. The work done against
        it becomes electrical energy, so more current makes the coil harder to turn.
\end{itemize}

%% ------------------------------------------------------------------
\section{AC induction motors and magnetic braking}\label{ch5:sec:induction}
\covers{M6.IQ4.02, M6.IQ4.03}

\begin{bigidea}
When a field and metal move past each other, \term{eddy currents} (loops of current
inside the metal) are induced. By Lenz's law, their forces oppose the \emph{relative}
motion.
\end{bigidea}

\subsection{AC induction motors}
In an \term{AC induction motor} the rotor current is induced, so there are no brushes
and no commutator.
\begin{enumerate}
  \item The stator has three pairs of coils on \term{three-phase AC} (three AC supplies
        a third of a cycle apart). They peak in turn, so they make a \term{rotating
        magnetic field}: fixed strength, direction turning steadily. No coil moves.
  \item The rotor is a \term{squirrel cage}: metal bars joined at both ends by end
        rings. The moving field changes the flux through each loop of bars, so currents
        are induced (Faraday).
  \item These currents are in the field, so the bars feel forces ($F = lIB$). By Lenz's
        law the forces oppose the \emph{relative} motion: the rotor is dragged after
        the field.
\end{enumerate}
\begin{itemize}
  \item \term{Slip}: the rotor always turns a little slower than the field. If it caught
        up, there would be no flux change, no current and no force. A bigger load gives
        more slip, so more current and more torque.
  \item No brushes or commutator, so little wear and no sparking. The supply frequency
        sets the speed, so it is hard to change.
\end{itemize}

\subsection{Magnetic braking}
\begin{sidebyside}[0.45]
  \centering
  \begin{tikzpicture}[x=1.0cm, y=1.0cm]
    \filldraw[fill=hscmuted!12, draw=hscink, line width=0.9pt] (0,0) circle (1.9);
    \fill[hscink] (0,0) circle (1.5pt);
    \draw[->, hscmuted, line width=0.9pt] (125:2.15) arc (125:55:2.15);
    \node[font=\footnotesize, text=hscmuted, anchor=south] at (90:2.17) {disc turns clockwise};
    \filldraw[fill=hscfield!14, draw=hscfield, dashed, line width=0.6pt] (0.35,-0.75) rectangle (1.75,0.75);
    \foreach \xx in {0.55,1.55} \foreach \yy in {-0.42,0.42}
      \node[text=hscfield, font=\footnotesize] at (\xx,\yy) {$\times$};
    \node[font=\scriptsize, text=hscfield, anchor=east, align=right] at (0.3,0.55) {$B$ into\\the page};
    \draw[hscseriesB, line width=1.2pt, ->] (1.05,0.75) ++(-90:0.3) arc (-90:250:0.3);
    \draw[hscseriesA, line width=1.2pt, ->] (1.05,-0.75) ++(90:0.3) arc (90:-250:0.3);
    \draw[->, hscink, line width=0.9pt] (0.5,-0.15) -- (1.6,-0.15) node[pos=0.25, above, font=\scriptsize] {$I$};
    \draw[force] (1.05,-0.08) -- (1.05,0.4) node[pos=0.55, right, text=hscforce, font=\footnotesize] {$F$};
    \draw[velocity] (1.85,0.45) -- (1.85,-0.45) node[pos=0.5, right, text=hscvel, font=\footnotesize] {$v$};
    \node[font=\scriptsize, text=hscseriesB, anchor=west] at (1.42,1.28) {metal enters field};
    \node[font=\scriptsize, text=hscseriesA, anchor=west] at (1.42,-1.28) {metal leaves field};
  \end{tikzpicture}
  \captionof{figure}{A metal disc turning through a field into the page ($\times$) on its
    right side.}\label{ch5:fig:disc}
\nextside
  \term{Magnetic braking} slows moving metal with eddy currents, without touching it
  (\cref{ch5:fig:disc}).
  \begin{itemize}
    \item Metal \emph{entering} the field: flux into the page grows, so the eddy current
          is anticlockwise. Metal \emph{leaving}: clockwise.
    \item In the field both currents flow to the right. The force on them points up,
          opposite to $v$, so the disc slows.
    \item Energy: kinetic $\to$ electrical $\to$ heat in the metal.
  \end{itemize}
\end{sidebyside}

\begin{itemize}
  \item The braking force grows with speed. At zero speed there is no flux change and no
        force, so it cannot hold an object still. Friction brakes finish the stop.
  \item A stronger field or a better conductor brakes more strongly. Slots break up the
        eddy-current loops and weaken it. No contact, so no wear.
  \item Heat $Q$ on a level track $=$ kinetic energy lost:
        $Q = \tfrac{1}{2}mu^2 - \tfrac{1}{2}mv^2$ ($u$ = speed before braking, $v$ = speed
        after). If it also falls a height $\Delta h$, add $mg\Delta h$.
  \item \term{Regenerative braking}: the motor of an electric car acts as a generator
        and charges the battery. By Lenz's law the force opposes the motion, so the car
        slows, and the energy is saved.
\end{itemize}

%% ------------------------------------------------------------------
\section{Key experiments}\label{ch5:sec:experiments}
\covers{M6.IQ4.01, M6.IQ4.03}

Two practicals matter here. In the \textbf{magnetic braking practical}, a magnet falls
slowly through a copper tube. A swinging aluminium plate between magnet poles stops
quickly, but a slotted plate swings longer. Eddy currents oppose the motion, and the
energy becomes heat. The \textbf{DC motor practical} is below.

\begin{experiment}[unbreakable]{Investigating a simple DC motor}
\aim{To find the job of each part of a DC motor, and the effect of back emf.}
\method{A coil of enamelled wire (wire with a thin insulating coat) rests on paper-clip
  supports above a magnet. The coat is scraped off all round one end, but off only the
  top half of the other end. Then a lab motor's current is read while it runs, and
  again while it is held still for 1--2 seconds.}
\results{The coil spins only with one end half scraped. Lab motor: held still
  \qty{1.5}{A}; running \qty{0.30}{A}.}
\conclusion{The half-scraped end switches the current off for the half turn when the
  torque would act backwards. So, like a commutator, it keeps the coil turning one way.
  The clips act as brushes. The lab motor draws less current when it runs, because its
  back emf opposes the supply.}
\significance{It shows every part of a DC motor, and back emf, in a few minutes.}
\end{experiment}

%% ------------------------------------------------------------------
\begin{summarybox}{Remember}
\begin{itemize}
  \item $\tau = nIAB\sin\theta$, $\theta$ from the normal: largest with the plane $\parallel B$.
  \item The split ring reverses the current every half turn, so the torque always turns
        the coil one way.
  \item Back emf opposes the supply and grows with speed:
        $I = \frac{V - \varepsilon_\mathrm{back}}{R}$ (not on the sheet).
  \item Generator emf is largest with the plane $\parallel B$. Slip rings: AC; split ring: DC.
  \item Induction motor: a rotating field drags the cage rotor round, a little slower (slip).
  \item Magnetic brake: eddy currents oppose the motion; no force at zero speed.
  \item Lenz's law is conservation of energy: if induced currents helped the change,
        energy would be made from nothing.
\end{itemize}
\end{summarybox}
}{}% Motors, Generators and Braking

%% ---------------------------------------------------------------- back
\IfFileExists{./formulas.tex}{% formulas.tex -- Module 6 ESSENTIALS: All formulas at a glance.
% One compact table of every formula in Module 6, in chapter order. The rows, their
% notation and their sheet status are copied from the final full book
% (src/mod6/formulasheet.tex: every formula card, then the other formulas that the
% Left out (not used in the Essentials text): F = GMm/r^2 and r^3/T^2 (gravity, used in the
% full book only for a comparison), P = Fv, I_perp = I sin(theta).
% Notation note (as in the full book): motional emf is written epsilon = lvB;
\unnumberedchapter[\faIcon{square-root-alt}]{All formulas at a glance}
\phantomsection\label{ess:fs:chap}

% ---- helpers for this chapter only (prefix essfs)
\newcommand\essfsYes{\textcolor{hscformula}{\footnotesize\bfseries\faCheck\ Yes}}
\newcommand\essfsNo{\textcolor{hscmistake}{\footnotesize\bfseries\faTimes\ Learn it}}
\newcommand\essfsComb{\textcolor{hsctiptext}{\footnotesize\bfseries\faLink\ Combine}}
% a little space above and below each formula, so tall fractions never touch the rules
\newcommand\essfspad[1]{\raisebox{0pt}[\dimexpr\height+0.9ex\relax][\dimexpr\depth+0.7ex\relax]{#1}}
% \essfsrow{mark}{formula}{what it is for}{chapter number}
\newcommand\essfsrow[4]{\rule[-1.3ex]{0pt}{4.0ex}\essfspad{#2} & #3 & #1 &
  \hyperref[ch#4:chap]{\textcolor{hscbrand}{\bfseries #4}}\\\hline}
% a navy rule between two chapters (drawn over the last grey rule)
\newcommand\essfsnext{\noalign{\vskip-\arrayrulewidth}\arrayrulecolor{hscbrand}%
  \specialrule{0.7pt}{0pt}{0pt}\arrayrulecolor{hscrule}}
\newcommand\essfshead{\colhead{Formula} & \colhead{What it is for} &
  \colhead{On the sheet} & \colhead{Ch}}

Every formula used in these Essentials, in chapter order. Use SI units in every formula.
``Ch'' is the chapter (click the number to go there). ``On the sheet'' is about the
official HSC formula sheet: \essfsYes\ = printed on the sheet; \essfsNo\ = not on the
sheet, so learn it or how to derive it; \essfsComb\ = not printed in this form, so join
formulas from the sheet.
\par\smallskip

{\small
\setlength\LTpre{9pt}\setlength\LTpost{6pt}
\renewcommand\arraystretch{1.0}\arrayrulecolor{hscrule}
\begin{longtable}{@{}C{47mm}@{\hspace{3mm}}L{\dimexpr\linewidth-47mm-17mm-5mm-9mm\relax}%
    @{\hspace{3mm}}C{17mm}@{\hspace{3mm}}C{5mm}@{}}
  \essfshead\\\hline\endfirsthead
  \essfshead\\\hline\endhead
  % ------------------------------------------------------------------ Chapter 1
  \essfsrow{\essfsYes}{$E = \dfrac{V}{d}$}{uniform field between parallel plates}{1}
  \essfsrow{\essfsYes}{$\vec{F} = q\vec{E}$ \quad $\vec{F}_\mathrm{net} = m\vec{a}$}
    {force on a charge; its acceleration $a = qE/m$}{1}
  \essfsrow{\essfsNo}{$qE = mg$ \quad $q = ne$}{a charged drop held still}{1}
  \essfsrow{\essfsYes}{$V = \dfrac{\Delta U}{q}$}{voltage: energy change per coulomb}{1}
  \essfsrow{\essfsYes}{$W = qV$ \quad $W = qEd$}
    {work done by the field ($qEd$: uniform only)}{1}
  \essfsrow{\essfsYes}{$K = \dfrac{1}{2}mv^2$}{kinetic energy of a charge}{1}
  \essfsrow{\essfsComb}{$\dfrac{1}{2}mv^2 = qV$ \quad $v = \sqrt{\dfrac{2qV}{m}}$}
    {speed gained \emph{from rest} through $V$}{1}
  \essfsrow{\essfsComb}{$qV = \dfrac{1}{2}mv^2 - \dfrac{1}{2}mu^2$}
    {a charge that is already moving (speed $u$)}{1}
  \essfsrow{\essfsYes}{\begin{tabular}{@{}c@{}}$v = u + at$ \quad $s = ut + \dfrac{1}{2}at^2$\\\noalign{\vskip 4pt}
    $v^2 = u^2 + 2as$\end{tabular}}
    {constant acceleration across the plates}{1}
  \essfsrow{\essfsComb}{$t = \dfrac{L}{u}$ \quad $y = \dfrac{1}{2}at^2 = \dfrac{qEL^2}{2mu^2}$}
    {time between plates of length $L$; deflection $y$}{1}
  \essfsrow{\essfsNo}{$y = \left(\dfrac{qE}{2mu^2}\right)x^2$}
    {path between the plates: a parabola}{1}
  \essfsrow{\essfsNo}{\begin{tabular}{@{}c@{}}$u_x = u\cos\theta$ \quad $u_y = u\sin\theta$\\\noalign{\vskip 4pt}
    $v = \sqrt{v_x^2 + v_y^2}$ \quad $\theta = \tan^{-1}\!\left(\dfrac{v_y}{v_x}\right)$\end{tabular}}
    {split a velocity into parts; join them again}{1}
  \essfsrow{\essfsYes}{$F = qv_\perp B = qvB\sin\theta$}
    {magnetic force on a moving charge}{1}
  \essfsrow{\essfsYes}{$W = F_\parallel s = Fs\cos\theta$}
    {magnetic force does no work: speed constant}{1}
  \essfsrow{\essfsYes}{$F_\mathrm{c} = \dfrac{mv^2}{r}$ \quad $v = \dfrac{2\pi r}{T}$}
    {circle: $qvB$ is the centripetal force}{1}
  \essfsrow{\essfsNo}{$r = \dfrac{mv}{qB}$ \quad $T = \dfrac{2\pi m}{qB}$}
    {radius and period of the circle in $B$}{1}
  \essfsrow{\essfsNo}{$qE = qvB$ \quad $v = \dfrac{E}{B}$}
    {crossed fields: the velocity selector}{1}
  \essfsnext
  % ------------------------------------------------------------------ Chapter 2
  \essfsrow{\essfsYes}{$B = \dfrac{\mu_0 I}{2\pi r}$}{field near a long straight wire}{2}
  \essfsrow{\essfsYes}{$B = \dfrac{\mu_0 NI}{L}$}
    {field inside a solenoid ($L$ = its length)}{2}
  \essfsrow{\essfsYes}{$F = lI_\perp B = lIB\sin\theta$}
    {force on a current-carrying wire}{2}
  \essfsrow{\essfsYes}{$\dfrac{F}{l} = \dfrac{\mu_0}{2\pi}\dfrac{I_1 I_2}{r}$}
    {force per metre between parallel wires}{2}
  \essfsrow{\essfsYes}{$I = \dfrac{q}{t}$}{current is charge per second}{2}
  \essfsrow{\essfsComb}{$F = \Delta m\, g$}
    {current balance: reading ($\Delta m$, kg) to force}{2}
  \essfsnext
  % ------------------------------------------------------------------ Chapter 3
  \essfsrow{\essfsYes}{$\Phi = B_\parallel A = BA\cos\theta$}
    {flux through a loop ($\theta$ from the \emph{normal})}{3}
  \essfsrow{\essfsNo}{$\Delta\Phi = \Phi_\mathrm{final} - \Phi_\mathrm{initial}$}
    {change in flux (reversed: $\lvert\Delta\Phi\rvert = 2BA$)}{3}
  \essfsrow{\essfsYes}{$\varepsilon = -N\dfrac{\Delta\Phi}{\Delta t}$}
    {Faraday's law (the minus sign is Lenz's law)}{3}
  \essfsrow{\essfsYes}{$V = IR$}{Ohm's law: a resistor, a coil or a line}{3}
  \essfsrow{\essfsComb}{$I = \dfrac{\varepsilon}{R}$}{induced current in a closed circuit}{3}
  \essfsrow{\essfsNo}{$\varepsilon = lvB$}
    {emf of a rod moving across a field}{3}
  \essfsrow{\essfsComb}{$\lvert\varepsilon_\mathrm{s}\rvert = N_\mathrm{s}\,
      \dfrac{\mu_0 N_\mathrm{p} A}{L}\,\dfrac{\lvert\Delta I_\mathrm{p}\rvert}{\Delta t}$}
    {emf in a coil wound on a solenoid}{3}
  \essfsnext
  % ------------------------------------------------------------------ Chapter 4
  \essfsrow{\essfsYes}{$\dfrac{V_\mathrm{p}}{V_\mathrm{s}} = \dfrac{N_\mathrm{p}}{N_\mathrm{s}}$}
    {transformer (turns) equation, on AC}{4}
  \essfsrow{\essfsYes}{$P = VI$}{electrical power; current sent: $I = P/V$}{4}
  \essfsrow{\essfsYes}{$V_\mathrm{p} I_\mathrm{p} = V_\mathrm{s} I_\mathrm{s}$}
    {ideal transformer: power in = power out}{4}
  \essfsrow{\essfsComb}{$\dfrac{I_\mathrm{s}}{I_\mathrm{p}} = \dfrac{V_\mathrm{p}}{V_\mathrm{s}}
    = \dfrac{N_\mathrm{p}}{N_\mathrm{s}}$}
    {currents in an ideal transformer}{4}
  \essfsrow{\essfsNo}{\begin{tabular}{@{}c@{}}$\eta = \dfrac{P_\mathrm{out}}{P_\mathrm{in}} \times 100\%$\\\noalign{\vskip 8pt}
      $\phantom{\eta} = \dfrac{V_\mathrm{s}I_\mathrm{s}}{V_\mathrm{p}I_\mathrm{p}} \times 100\%$\end{tabular}}
    {efficiency of a real transformer}{4}
  \essfsrow{\essfsYes}{$P = \dfrac{\Delta E}{\Delta t}$}{energy used or wasted in a time}{4}
  \essfsrow{\essfsYes}{$V_\mathrm{drop} = IR$}
    {voltage lost along a power line}{4}
  \essfsrow{\essfsNo}{$P_\mathrm{loss} = I^2R$}{power lost as heat in a line}{4}
  \essfsrow{\essfsYes}{$f = \dfrac{1}{T}$}{AC frequency (mains: \qty{50}{Hz})}{4}
  \essfsnext
  % ------------------------------------------------------------------ Chapter 5
  \essfsrow{\essfsYes}{$\tau = nIA_\perp B = nIAB\sin\theta$}
    {torque on a coil ($\theta$ from the normal)}{5}
  \essfsrow{\essfsYes}{$\tau = r_\perp F = rF\sin\theta$}{torque from one force}{5}
  \essfsrow{\essfsNo}{$I = \dfrac{V - \varepsilon_\mathrm{back}}{R}$}
    {current in a running DC motor (back emf)}{5}
  \essfsrow{\essfsNo}{$VI = \varepsilon_\mathrm{back}I + I^2R$}
    {motor: power in = mechanical + heat}{5}
  \essfsrow{\essfsYes}{$\varepsilon = -N\dfrac{\Delta\Phi}{\Delta t}$}
    {average emf of a turning generator coil}{5}
  \essfsrow{\essfsComb}{$Q = \dfrac{1}{2}mu^2 - \dfrac{1}{2}mv^2$}
    {heat from a magnetic brake (level track)}{5}
  \essfsrow{\essfsComb}{$Q = \left(\dfrac{1}{2}mu^2 - \dfrac{1}{2}mv^2\right) + mg\Delta h$}
    {magnetic brake on a slope or drop tower}{5}
\end{longtable}}

\section*{Data sheet values used in this module}

\begin{table}
  \centering\small
  \begin{tabular}{@{}L{55mm}L{44mm}L{41mm}@{}}
    \toprule
    \colhead{Quantity (as on the sheet)} & \colhead{Value} & \colhead{Used in}\\
    \midrule
    Charge on electron & $q_\mathrm{e} = \qty{-1.602e-19}{C}$ & \cref{ch1:chap}\\
    Mass of electron & $m_\mathrm{e} = \qty{9.109e-31}{kg}$ & \cref{ch1:chap}\\
    Mass of proton & $m_\mathrm{p} = \qty{1.673e-27}{kg}$ & \cref{ch1:chap}\\
    Earth's gravitational acceleration & $g = \qty{9.8}{m.s^{-2}}$
      & \cref{ch1:chap,ch2:chap,ch5:chap}\\
    Magnetic permeability constant & $\mu_0 = 4\pi \times 10^{-7}$ \si{N.A^{-2}}
      & \cref{ch2:chap,ch3:chap}\\
    \qty{1}{eV} & \qty{1.602e-19}{J} & \cref{ch1:chap}\\
    \bottomrule
  \end{tabular}
\end{table}

\begin{itemize}
  \item The sheet prints the electron's charge as negative. In the formulas, use its size,
        $e = \qty{1.602e-19}{C}$. Find the direction with a rule.
  \item A proton has charge $+e$. An alpha particle has charge $+2e$.
\end{itemize}
}{}
\IfFileExists{./experiments.tex}{% experiments.tex -- Module 6 ESSENTIALS: Key experiments at a glance.
% One row per Key experiment box of ch1.tex ... ch5.tex in this folder (the Chapter 1
% box has two parts, so it has two rows). Facts and numbers as in those boxes, which
% follow the full book (src/mod6/ch*.tex and the overview table in
% src/mod6/frontmatter.tex).
\unnumberedchapter[\faFlask]{Key experiments at a glance}

The syllabus asks you to carry out or study these investigations. Each one has a Key
experiment box (or a short paragraph) in its chapter (``Ch'': click the number to go
there).

\par\smallskip
{\small
\setlength\LTpre{4pt}\setlength\LTpost{4pt}
\renewcommand\arraystretch{1.25}\arrayrulecolor{hscrule}
\newcommand\essexhead{\colhead{Experiment} & \colhead{What was done} &
  \colhead{What it showed} & \colhead{Ch}}
\newcommand\essexch[1]{\hyperref[ch#1:sec:experiments]{\textcolor{hscbrand}{\bfseries #1}}}
\begin{longtable}{@{}L{31mm}@{\hspace{3mm}}L{50mm}@{\hspace{3mm}}%
    L{\dimexpr\linewidth-31mm-50mm-8mm-9mm\relax}@{\hspace{3mm}}C{8mm}@{}}
  \essexhead\\\hline\endfirsthead
  \essexhead\\\hline\endhead
  \textbf{Charges between parallel plates} (simulation)
    & An electron enters between the plates, moving parallel to them. Change $V$, the
      speed and the particle; record the deflection $y$.
    & $y \propto V$ and $y \propto x^2$: the path is a parabola, like a projectile.
      Half the speed gives four times the deflection.
    & \essexch{1}\\\hline
  \textbf{Speed gained from a voltage} (simulation)
    & An electron gun: record the exit speed $v$ for several voltages $V$, then double
      the gap.
    & $v^2 \propto V$, as $qV = \frac{1}{2}mv^2$ predicts. The gap makes no difference.
    & \essexch{1}\\\hline
  \textbf{The force on a wire} (current balance)
    & A wire passes through the gap of a magnet that sits on a balance. Change $I$, then
      the angle; $F = \Delta m\, g$.
    & $F \propto I$. $F$ is largest at \ang{90} to $B$ and zero when parallel. This
      agrees with $F = lIB\sin\theta$.
    & \essexch{2}\\\hline
  \textbf{Two parallel wires} (\textbf{required})
    & Two rods, one on a balance and one clamped above it. Change $I$ (and the gap $r$);
      plot $F$ against $I^2$.
    & Like currents attract, opposite currents repel. $F \propto I_1 I_2$, as
      $F/l = (\mu_0/2\pi)(I_1 I_2/r)$ predicts.
    & \essexch{2}\\\hline
  \textbf{Faraday's induction experiments} (1831)
    & A magnet moved in and out of a coil with a galvanometer (a meter for very small
      currents); a switch closed and opened on two coils on an iron ring.
    & An emf is induced only while the flux changes. A faster change gives a bigger
      emf; the opposite change, the opposite current.
    & \essexch{3}\\\hline
  \textbf{Falling magnets and swinging plates} (eddy currents)
    & A magnet falls through PVC, copper and slotted copper tubes. Solid and slotted
      plates swing between magnet poles.
    & Eddy currents oppose the motion (Lenz's law): the fall in copper is much slower.
      The energy becomes heat. Slots weaken the braking.
    & \essexch{3},\newline\essexch{5}\\\hline
  \textbf{Testing the transformer equations}
    & A demountable (take-apart) transformer: measure $V_\mathrm{s}$ for several $N_\mathrm{s}$;
      read all four meters with a load; open the core.
    & $V_\mathrm{s} \propto N_\mathrm{s}$, a little below the ideal line. Efficiency
      about 90\%. An open core gives a much smaller $V_\mathrm{s}$ (flux leakage).
    & \essexch{4}\\\hline
  \textbf{Investigating a simple DC motor}
    & A coil on paper-clip supports above a magnet; one wire end is only half scraped.
      A lab motor's current is read while it runs and while it is held still.
    & The half-scraped end switches the current off for half of each turn, so the torque
      always acts one way (like a commutator); the clips act as brushes. Back emf makes
      the running current smaller than the stalled current.
    & \essexch{5}\\\hline
\end{longtable}}
}{}
\IfFileExists{./glossary.tex}{% glossary.tex -- Module 6 ESSENTIALS glossary (A-Z). GENERATED by
% src/mod6/glossary.tex (the key terms only; meanings and Chinese unchanged,
% a few meanings shortened). Edit the script, not this file.
\unnumberedchapter[\faBook]{Glossary}

The key words of this module, from A to Z. The meaning is in plain English. The last
column gives the usual Chinese term used in physics textbooks. The full guide has a
longer glossary.

{\small
\begin{glossarytable}
  % ---- A
  \gterm{AC induction motor}{a motor whose rotor currents are induced by a rotating field}{交流感应电动机（异步电动机）}
  \gterm{Alternating current (AC)}{current that keeps reversing its direction}{交流电}
  \gterm{Ampere (A)}{the SI unit of current: \qty{1}{A} means one coulomb of charge passes a point each second}{安培}
  % ---- B
  \gterm{Back emf ($\varepsilon_\mathrm{back}$)}{the emf induced in a spinning motor coil; it opposes the supply}{反电动势}
  \gterm{Brushes}{fixed contacts that join the spinning coil to the circuit}{电刷}
  % ---- C
  \gterm{Centripetal force}{the net force towards the centre of a circle}{向心力}
  \gterm{Charge-to-mass ratio ($q/m$)}{charge divided by mass; a bigger $q/m$ gives a bigger acceleration}{荷质比（比荷）}
  \gterm{Conventional current}{the way positive charge would flow: from $+$ round the circuit to $-$}{（规定的）电流方向}
  \gterm{Core (soft-iron core)}{the iron that carries the magnetic field from one coil to the other}{铁芯（软铁芯）}
  \gterm{Couple}{two equal and opposite forces that turn an object without moving it sideways}{力偶}
  % ---- D
  \gterm{DC motor}{a device that turns electrical energy from a DC supply into rotational kinetic energy}{直流电动机}
  \gterm{Deflection}{how far a particle is pushed sideways from its straight path}{偏转（偏移量）}
  \gterm{Direct current (DC)}{current that flows in one direction only}{直流电}
  % ---- E
  \gterm{Eddy currents}{loops of current induced inside solid metal when the flux through it changes}{涡流}
  \gterm{Effective value (of AC)}{the steady (DC) voltage that would heat a heater by the same amount}{有效值}
  \gterm{Efficiency ($\eta$)}{the percentage of the input power that comes out as useful power}{效率}
  \gterm{Electric field}{a region where a charge feels an electric force}{电场}
  \gterm{Electric field strength ($E$)}{the force on each coulomb of charge; unit \si{V.m^{-1}}, the same as \si{N.C^{-1}}}{电场强度}
  \gterm{Electric potential energy ($U$)}{energy a charge stores because of its position in an electric field}{电势能}
  \gterm{Electromagnetic induction}{making an emf by changing the magnetic flux through a circuit}{电磁感应}
  \gterm{Electromotive force (emf, $\varepsilon$)}{not a force: the energy a source gives each coulomb of charge; unit volt (V)}{电动势}
  \gterm{Electron volt (eV)}{the energy one electron gains through \qty{1}{V}: \qty{1}{eV} $=$ \qty{1.602e-19}{J}}{电子伏特}
  \gterm{Elementary charge ($e$)}{the size of the charge on one proton or electron, \qty{1.602e-19}{C}}{元电荷}
  % ---- F
  \gterm{Faraday's law}{the induced emf equals the number of turns times the rate of change of flux}{法拉第电磁感应定律}
  % ---- G
  \gterm{Generator}{a device that turns kinetic energy into electrical energy}{发电机}
  % ---- I
  \gterm{Ideal transformer}{a perfect transformer that wastes no energy, so power in equals power out}{理想变压器}
  \gterm{Incomplete flux linkage (flux leakage)}{some of the flux from the primary misses the secondary}{漏磁（磁通未完全耦合）}
  \gterm{Induced current}{the current that an induced emf drives around a closed circuit}{感应电流}
  \gterm{Induced emf}{an emf produced by a changing magnetic flux}{感应电动势}
  % ---- L
  \gterm{Laminated core}{a core made of thin iron sheets with insulation between them}{叠片铁芯（硅钢片铁芯）}
  \gterm{Lenz's law}{an induced current's magnetic field opposes the change in flux that made it}{楞次定律}
  \gterm{Lever arm}{the perpendicular distance from the axis to the line of a force}{力臂}
  % ---- M
  \gterm{Magnetic braking}{slowing moving metal with eddy currents, without touching it}{电磁制动（电磁阻尼）}
  \gterm{Magnetic field}{a region where a magnet or a moving charge feels a magnetic force}{磁场}
  \gterm{Magnetic field strength ($B$)}{how strong a magnetic field is; also called magnetic flux density; unit tesla (T)}{磁感应强度}
  \gterm{Magnetic flux ($\Phi$)}{how much magnetic field passes through an area; unit weber (Wb)}{磁通量}
  \gterm{Magnetic flux density}{another name for $B$: the flux through each square metre at right angles to the field}{磁通密度（即磁感应强度）}
  \gterm{Magnetic permeability constant ($\mu_0$)}{the constant that links a current to the field it makes in a vacuum, $4\pi \times 10^{-7}$ \si{N.A^{-2}}}{真空磁导率}
  \gterm{Mass spectrometer}{a machine that sorts ions by mass, using the radius of their path in a magnetic field}{质谱仪}
  \gterm{Motional emf}{the emf in a conductor because it moves across a magnetic field}{动生电动势}
  \gterm{Motor effect}{the force on a current-carrying wire in an external magnetic field}{磁场对电流的作用（安培力）}
  \gterm{Mutual induction}{an emf induced in one coil by a changing current in a nearby coil}{互感}
  % ---- N
  \gterm{Normal}{an imaginary line at \ang{90} to a flat surface, such as the plane of a coil}{法线}
  % ---- P
  \gterm{Parabola}{the curved path of a thrown ball, made by a constant force in one direction}{抛物线}
  \gterm{Perpendicular}{at right angles (\ang{90})}{垂直}
  \gterm{Primary coil}{the input coil, joined to the supply}{原线圈（初级线圈）}
  % ---- R
  \gterm{Radial magnetic field}{a field whose lines point along radii, like the spokes of a wheel}{辐向磁场}
  \gterm{Regenerative braking}{braking in which the motor acts as a generator and the battery stores the energy}{再生制动（能量回收制动）}
  \gterm{Resistive heating}{heat made when a current flows through a resistance ($I^2R$)}{电阻发热（焦耳热）}
  \gterm{Right-hand grip rule}{thumb along the current; the curled fingers show the field around it}{右手螺旋定则（安培定则）}
  \gterm{Right-hand palm rule}{fingers along $B$, thumb along $I$ (or $v$ of a positive charge): the palm pushes along $F$}{右手掌定则（结果同左手定则）}
  \gterm{Rotating magnetic field}{a field of fixed strength whose direction turns steadily, made by coils that do not move}{旋转磁场}
  \gterm{Rotor (armature)}{the part of a motor that turns: a coil on an iron core, fixed to the axle}{转子（电枢）}
  % ---- S
  \gterm{Secondary coil}{the output coil, in which the emf is induced}{副线圈（次级线圈）}
  \gterm{Slip}{the difference between the speed of the rotating field and the speed of the rotor}{转差}
  \gterm{Slip rings}{two full rings, each always joined to the same end of a generator coil}{滑环}
  \gterm{Solenoid}{a long coil of wire with many turns; the field inside is uniform}{螺线管}
  \gterm{Split-ring commutator}{two half-rings that reverse the current in the coil every half turn}{换向器（裂环换向器）}
  \gterm{Squirrel-cage rotor}{metal bars joined at both ends by end rings, like a cage}{鼠笼式转子}
  \gterm{Stator}{the part of a motor that does not move; it makes the magnetic field}{定子}
  \gterm{Step-down transformer}{a transformer with fewer secondary turns, so it makes the voltage smaller}{降压变压器}
  \gterm{Step-up transformer}{a transformer with more secondary turns, so it makes the voltage bigger}{升压变压器}
  \gterm{Substation}{a place with transformers that change the voltage of the grid}{变电站}
  % ---- T
  \gterm{Tesla (T)}{the SI unit of magnetic field strength: \qty{1}{T} $=$ \qty{1}{N.A^{-1}.m^{-1}} $=$ \qty{1}{Wb.m^{-2}}}{特斯拉}
  \gterm{Three-phase AC}{three AC supplies that peak one after another, \ang{120} apart}{三相交流电}
  \gterm{Torque ($\tau$)}{the turning effect of a force; unit \si{N.m}}{力矩}
  \gterm{Transformer}{a device that changes an AC voltage}{变压器}
  \gterm{Transmission line}{a long, thick wire on tall towers that carries electricity a long way}{输电线}
  \gterm{Turns ratio}{$N_\mathrm{p} : N_\mathrm{s}$, the primary turns compared with the secondary turns}{匝数比}
  % ---- U
  \gterm{Uniform field}{a field with the same strength and direction everywhere}{匀强场（匀强电场、匀强磁场）}
  % ---- V
  \gterm{Velocity selector}{crossed fields that let only one speed pass straight through}{速度选择器}
  \gterm{Voltage drop}{the voltage used up pushing a current through a resistance, $IR$}{电压降（电压损失）}
  % ---- W
  \gterm{Weber (Wb)}{the SI unit of magnetic flux: \qty{1}{Wb} $=$ \qty{1}{T.m^2}}{韦伯}
\end{glossarytable}}
}{}

\end{document}
